\documentclass[11pt,a4paper]{article}

\usepackage[utf8]{inputenc}
\usepackage[T1]{fontenc}
\usepackage{lmodern}
\usepackage[margin=2.4cm]{geometry}
\usepackage{amsmath,amssymb,amsthm,mathtools}
\usepackage{bm}
\usepackage{booktabs,array,float,multirow}
\usepackage{enumitem}
\usepackage{microtype}
\usepackage{syntax}
\usepackage{graphicx}
\usepackage{xcolor}
\usepackage[numbers,sort&compress]{natbib}
\usepackage{listings}
\usepackage[colorlinks=true,linkcolor=blue!45!black,%
            citecolor=blue!45!black,urlcolor=blue!45!black]{hyperref}
\usepackage{cleveref}

% ─── Theorem environments ────────────────────────────────────────────
\newtheoremstyle{thmstyle}{\topsep}{\topsep}{\itshape}{}{\bfseries}{.}{.5em}{}
\theoremstyle{thmstyle}
\newtheorem{theorem}{Theorem}[section]
\newtheorem{lemma}[theorem]{Lemma}
\newtheorem{proposition}[theorem]{Proposition}
\newtheorem{corollary}[theorem]{Corollary}
\theoremstyle{definition}
\newtheorem{definition}[theorem]{Definition}
\newtheorem{axiom}{Axiom}
\newtheorem{example}[theorem]{Example}
\newtheorem{principle}[theorem]{Principle}
\theoremstyle{remark}
\newtheorem{remark}[theorem]{Remark}

\crefname{axiom}{Axiom}{Axioms}
\Crefname{axiom}{Axiom}{Axioms}
\crefname{principle}{Principle}{Principles}
\Crefname{principle}{Principle}{Principles}

% ─── Notation ────────────────────────────────────────────────────────
\newcommand{\R}{\mathbb{R}}
\newcommand{\N}{\mathbb{N}}
\newcommand{\Gph}{G}
\newcommand{\V}{V}
\newcommand{\E}{E}
\newcommand{\med}{\mathfrak{M}}
\newcommand{\cut}{\partial}
\newcommand{\sep}{\sigma}
\newcommand{\flo}{\beta}
\newcommand{\kap}{\kappa}
\newcommand{\Env}{\Gamma}
\newcommand{\store}{\Sigma}
\newcommand{\steps}{\longrightarrow}
\newcommand{\abs}[1]{\lvert#1\rvert}
\newcommand{\norm}[1]{\lVert#1\rVert}
\newcommand{\dd}{\,\mathrm{d}}
\newcommand{\Circ}{\textsf{Circuit}}
\newcommand{\Comp}{\textsf{Compartment}}
\newcommand{\Chg}{\textsf{Charge}}
\newcommand{\Obs}{\textsf{Observable}}
\newcommand{\Exp}{\textsf{Experiment}}
\newcommand{\Res}{\textsf{Result}}
\newcommand{\Strat}{\textsf{Stratum}}
\newcommand{\circuit}{\mathbf{circuit}}
\newcommand{\outb}{\mathbf{outbound}}
\newcommand{\retn}{\mathbf{return}}
\newcommand{\without}{\mathbf{without}}
\newcommand{\observe}{\mathbf{observe}}
\newcommand{\experiment}{\mathbf{experiment}}

% ─── Listings style for Vitruvius ───────────────────────────────────
\lstdefinelanguage{vitruvius}{
  morekeywords={circuit,outbound,return,compartment,capacitance,stratum,
    floor,experiment,observe,without,with,scaling,noise,across,let,
    element,conducts,at,delay,gain,type,event,against,report,
    intact,lesion,compare,run,for,module,import,where,derived,
    reflex,spinal,supraspinal,by,of,in,as,to},
  sensitive=true,
  morecomment=[l]{--},
  morestring=[b]",
}
\lstdefinestyle{vvs}{
  language=vitruvius,
  basicstyle=\ttfamily\small,
  keywordstyle=\color{blue!70!black}\bfseries,
  commentstyle=\color{green!45!black}\itshape,
  stringstyle=\color{red!55!black},
  numberstyle=\tiny\color{gray},
  frame=single,
  framesep=4pt,
  breaklines=true,
  columns=fullflexible,
  showstringspaces=false,
  xleftmargin=2pt,
}
\lstset{style=vvs}

% ─── Title ───────────────────────────────────────────────────────────
\title{\textbf{Vitruvius: A Language for Musculoskeletal Circuit
Experiments}\\[4pt]
\large Grammar, Type System, and Operational Semantics of \texttt{.vvs}}

\author{%
Kundai Farai Sachikonye\\
\small Technical University of Munich\\
\small \texttt{kundai.sachikonye@tum.de}}

\date{\today}

\begin{document}
\maketitle

\begin{abstract}
\noindent
We specify \emph{Vitruvius} (source extension \texttt{.vvs}), a domain-specific
language for expressing and running experiments on models of the human
musculoskeletal system. The language takes as its unit of description not the
signal but the \emph{circulation}: a closed path through the neuromuscular
graph comprising an efferent outbound phase and an afferent return phase. This
choice is forced rather than stylistic. We prove that a neuromuscular graph
admits no vertex connected to an unbounded charge reservoir, that a graph so
constrained possesses no static equilibrium under sustained metabolic input,
and consequently that its trajectories are bounded perpetual oscillations
rather than approaches to a set point. A language whose primitive is a signal
travelling from controller to plant cannot express this; a language whose
primitive is a circulation can. From the closure structure we derive three
static analyses that a compiler can perform before any numerical integration:
an \emph{aperture} analysis reporting which circulations lack a return phase
and are therefore predicted to fail to terminate rather than to degrade; a
\emph{compartment} analysis rejecting charge arithmetic that mixes distinct
capacitances, in the manner of dimensional analysis but over anatomical rather
than metric units; and a \emph{stratum} analysis enforcing that elements at
one characteristic time scale do not directly read the state of another,
implemented as an effect system whose deliberate violation is itself the model
of a coupling pathology. We give the complete lexical and context-free grammar,
a type system with seven types and four non-standard rules, and a small-step
operational semantics over a configuration comprising a circuit, an observation
store, and a monotone committed record. We prove progress and preservation, and
hence that a well-typed program does not get stuck; we prove that lesion
composition is associative, commutative, and idempotent, so that the lesioned
circuit of an experiment is independent of the order in which ablations are
declared; we prove that every experiment over a finite observable set
terminates; and we prove that the language's estimator for the catalytic power
of an event is non-degenerate, in contrast to the natural instance-specific
estimator, which we show is an algebraic identity that agrees with its own
prediction on every data set whatsoever and therefore has no null hypothesis.
We report the type separation statistic that determines when the
non-degenerate estimator has power, and specify the diagnostic that
distinguishes an uninformative model from a negative result. Three worked
programs---a stretch reflex, a postural loop under ablation, and a
cross-stratum coupling lesion---exercise every construct. A reference
implementation, exercised over ten programs yielding $257$ observations,
establishes that the language recovers clinical taxonomy from anatomy
alone---the upper/lower motor neuron distinction and the level-dependent
picture of spinal cord injury each follow from a membership test on declared
paths---and produces three statements not otherwise available: that regional
anaesthesia passes through a window in which contractile capacity is
preserved to full precision while the circulation is open, so coordination
fails without weakness; that surgical reinnervation restores closure through
a path $16$\,ms slower than the limb it replaces, so the repaired
circulation is not the original one; and that joint stiffness is an emergent
property of two coupled limit cycles rather than a commanded parameter.
\end{abstract}

\noindent\textbf{Keywords:} domain-specific language; musculoskeletal
simulation; motor control; closed-loop circuits; type systems; effect systems;
operational semantics; deafferentation; scientific modelling languages.

\tableofcontents

% =====================================================================
\section{Introduction}
\label{sec:intro}
% =====================================================================

\subsection{The design question}

Simulation platforms for human movement are organised around a pipeline. A
model is specified as a mechanical plant---segments, joints, musculotendon
actuators---and a controller supplies activations to that plant; the simulator
integrates the coupled system forward and returns trajectories
\citep{delp2007,seth2018,todorov2012}. The organising abstraction is the
\emph{signal}: activation flows from controller to plant, sensory information
flows back, and the two are distinct objects joined by a feedback connection
whose gain is a parameter.

This abstraction is adequate for a large class of questions and we do not
propose to replace it everywhere. It has, however, a specific and consequential
limitation, which is the subject of this paper. Consider the class of
experiments in which a structural element of the sensorimotor loop is removed
rather than attenuated: severing a proprioceptive afferent, abolishing a
return pathway, decoupling one anatomical level from another. In a
signal-and-gain architecture these manipulations are all expressed the same
way---by setting a gain to zero---and they therefore all predict the same kind
of outcome, namely a graded increase in variance about a preserved mean
trajectory. The empirical record does not support that prediction. Removal of
large-fibre proprioceptive input does not degrade coordinated movement; it
abolishes it, and does so while motor innervation and voluntary intent remain
intact \citep{cole1991,sainburg1995,ghez1995,lajoie1996,taub1965}.

A language in which severance and attenuation are the same operation cannot
express the distinction that the data demand. Vitruvius is built on a different
primitive for exactly this reason.

\subsection{The primitive}

The unit of description in Vitruvius is the \emph{circulation}: a closed path
through the neuromuscular graph with a designated outbound phase (efferent,
central to peripheral) and a designated return phase (afferent, peripheral to
central). A circulation is one object, not two signals joined by a connection.
Consequently there are two structurally distinct manipulations available to an
experimenter, and the language distinguishes them syntactically:

\begin{itemize}[leftmargin=1.4em,itemsep=2pt]
\item \textbf{Modifying} a circulation---changing a conductance, a delay, a
gain---leaves it closed and yields a circuit of the same topological class.
\item \textbf{Severing} a circulation---removing its return phase---opens it,
and yields a circuit of a different topological class.
\end{itemize}

The compiler can tell these apart without running anything, and
\Cref{sec:aperture} shows that it does: a program that opens a circulation
receives a static diagnostic naming the aperture and the prediction that
follows from it. This is the language's central design commitment and the
reason for its existence.

\subsection{Why closure is forced}

The commitment would be arbitrary if closure were merely a convenient way to
organise a model. \Cref{sec:foundations} shows it is not. We establish, from
premises that are individually standard, three results:

\begin{enumerate}[leftmargin=1.6em,itemsep=2pt]
\item The neuromuscular graph of an intact organism has no vertex connected to
an unbounded charge reservoir (\Cref{prop:no-ground}). Charge crossing any
membrane is redistributed within finite internal compartments; it does not
leave.
\item A graph so constrained, carrying voltage-gated conductances and supplied
with sustained metabolic energy, admits no static equilibrium
(\Cref{thm:no-equilibrium}), and its trajectories are therefore bounded and
non-static---perpetual oscillation within physiological limits
(\Cref{cor:perpetual}).
\item In such a graph a perturbation originating centrally does not terminate
peripherally unless the peripheral vertex lies on a closed path back to the
origin (\Cref{cor:closure-obligatory}).
\end{enumerate}

The third result is what distinguishes severance from attenuation. An
attenuated circulation is a closed path with small currents; a severed
circulation is not a path at all, and there is no value of a gain that
expresses the difference.

\subsection{What the language is not}

Vitruvius is not a general-purpose programming language. It has no
user-defined functions, no recursion, no higher-order values, and no
input/output beyond the declaration of models and the reporting of
observations. Numerical work---integration of membrane dynamics, evaluation of
cross-bridge kinetics, spectral estimation---is performed by a backend that
the language calls and does not inspect (\Cref{sec:backends}). In the sense in
which a build language coordinates compilation without compiling anything,
Vitruvius coordinates experiments without computing anything: it expresses
which circuits exist, which are ablated, and what is to be observed, and
delegates the rest.

Neither is it a replacement for existing musculoskeletal simulators. Its scope
is narrower and its commitments stronger. Where a general platform must
accommodate open-loop actuation, prescribed kinematics, and externally
specified controllers, Vitruvius admits none of these as well-formed
descriptions of an intact system, and gains in return the ability to
distinguish structural from parametric manipulation. We state the trade
plainly in \Cref{sec:scope}.

\subsection{Contributions}

\begin{enumerate}[leftmargin=1.6em,itemsep=1pt]
\item A derivation of circuit closure as a structural necessity of the intact
neuromuscular graph, from three standard premises (\Cref{sec:foundations}).
\item A complete lexical and context-free grammar for \texttt{.vvs}, proved
unambiguous (\Cref{sec:grammar}).
\item A type system with seven types and four non-standard rules: compartment
consistency, stratum containment, floor positivity, and typed-event
estimation (\Cref{sec:typing}).
\item The aperture analysis: a static decision procedure reporting the
circulations a program has opened, with the prediction that follows
(\Cref{sec:aperture}).
\item A small-step operational semantics with progress, preservation,
lesion-order independence, and termination (\Cref{sec:semantics}).
\item A proof that the natural estimator of an event's catalytic power is an
algebraic identity with no null hypothesis, and a typed estimator that is
non-degenerate, together with the statistic governing its power
(\Cref{sec:estimation}).
\item Three worked programs exercising every construct
(\Cref{sec:programs}).
\item A reference implementation whose fifty tests assert the specification's
properties directly, together with the empirical findings it produced and the
three specification-level errors that building it exposed
(\Cref{sec:implementation}).
\end{enumerate}

\subsection{Relation to existing work}

The formal apparatus is standard. The grammar is presented in the style of
\citet{aho2006}; the type system as inference rules with progress and
preservation in the manner of \citet{wright1994} and \citet{pierce2002}; the
dynamics as small-step structural operational semantics following
\citet{plotkin1981}. The compartment rule of \Cref{sec:compartments} is a
units-of-measure discipline in the sense of \citet{kennedy1997}, differing
only in that its units are anatomical rather than metric. The stratum rule of
\Cref{sec:strata} is an effect system in the sense of \citet{lucassen1988},
with the region structure of \citet{tofte1997}. The novelty claimed here is
not in the metatheory but in what the rules enforce and why the enforcement is
physiologically obligatory.

Among modelling languages, CellML \citep{hedley2001} and SBML
\citep{hucka2003} standardise the exchange of biological models but are
representations rather than languages with a type discipline; Modelica
\citep{fritzson2014} supplies acausal equation-based modelling with strong
unit checking, and is the closest relative of the compartment analysis
developed here, but has no notion of circuit closure and no analogue of the
aperture diagnostic. Musculoskeletal simulation platforms
\citep{delp2007,seth2018,todorov2012} supply the numerics that a Vitruvius
backend requires, and we regard them as complements rather than competitors:
\Cref{sec:backends} specifies the interface through which such a platform
discharges a Vitruvius model.

The motor-control theories with which the framework must be compared---the
equilibrium-point hypothesis \citep{feldman1986}, optimal feedback control
\citep{todorov2002}, and internal-model accounts
\citep{kawato1999,wolpert2000}---are discussed in \Cref{sec:discussion}, where
we identify the specific predictions on which they and the closure account
differ.

% =====================================================================
\section{Foundations}
\label{sec:foundations}
% =====================================================================

This section derives the structural facts the language encodes. Each premise
is standard; the conclusions are what motivate the design.

\subsection{The neuromuscular graph}

\begin{definition}[Neuromuscular graph]
\label{def:nm-graph}
The \emph{neuromuscular graph} is a finite directed graph
$\Gph=(\V,\E)$ whose vertices $\V$ are compartments at which membrane
potential is approximately uniform---a neuronal soma, a dendritic branch, an
axonal segment, a presynaptic terminal, a muscle fibre segment, a receptor
ending---and whose edges $\E$ are conductive pathways between compartments:
axonal conduction, dendritic coupling, synaptic contact, and the neuromuscular
junction.
\end{definition}

At each vertex $v$, conservation of current gives
\begin{equation}
\sum_{u:(u,v)\in\E} I_{uv}(t)
\;=\; C_{\mathrm{m}}^{(v)}\frac{\dd V^{(v)}}{\dd t}
\;+\; I_{\mathrm{ion}}^{(v)}(t),
\label{eq:kcl}
\end{equation}
where $C_{\mathrm{m}}^{(v)}$ is the membrane capacitance at $v$ and
$I_{\mathrm{ion}}^{(v)}$ the ionic current through channels at $v$. The ionic
currents follow the standard conductance-based description
\citep{hodgkin1952},
\begin{equation}
C_{\mathrm{m}}\frac{\dd V}{\dd t}
= -\sum_k g_k(V,t)\,\bigl(V-E_k\bigr) + I_{\mathrm{ext}},
\label{eq:hh}
\end{equation}
with $g_k$ the conductance for ion species $k$ and $E_k$ its reversal
potential. Specific membrane capacitance is approximately
$1\,\mu\mathrm{F/cm}^2$ across cell types
\citep{cole1941,gentet2000}, a fact we use in \Cref{sec:compartments}.

\subsection{Absence of an external ground}

\begin{definition}[External ground]
\label{def:ground}
A vertex $v^\ast$ of a circuit is an \emph{external ground} if it is connected
to an unbounded charge reservoir: a source or sink into which arbitrary
current may flow while its potential remains externally fixed.
\end{definition}

\begin{proposition}[The intact neuromuscular graph is ungrounded]
\label{prop:no-ground}
The \emph{in vivo} neuromuscular graph contains no external ground.
\end{proposition}

\begin{proof}
An external ground requires an unbounded reservoir. Every charge crossing a
membrane of the organism enters either the intracellular or the extracellular
fluid compartment, both of which are of finite volume and finite ionic
content. Ionic gradients are maintained by the metabolic work of active
transport, but a pump redistributes charge across a membrane rather than
supplying it from outside the organism: the ions it moves were already
present. Transdermal currents on the timescale of a movement are smaller than
intracellular currents by several orders of magnitude and cannot serve as an
unbounded sink. Finally, the reversal potentials $E_k$ appearing in
\eqref{eq:hh} are fixed by the Nernst relation from finite internal and
external concentrations, not by reference to an external potential. Hence no
vertex satisfies \Cref{def:ground}.
\end{proof}

\begin{remark}[The reference electrode is an artefact of measurement]
\label{rem:electrode}
Electrophysiological recording introduces a ground externally, and membrane
potentials are conventionally reported against it. This is a property of the
apparatus, not of the preparation. \Cref{prop:no-ground} concerns the intact
organism, in which no such reference exists. The distinction matters for the
language because a model that admits a grounded vertex is a model of a
recording configuration rather than of a behaving system, and Vitruvius
declines to typecheck one (\Cref{sec:typing}).
\end{remark}

\subsection{No static equilibrium}

\begin{axiom}[Sustained metabolic supply]
\label{ax:metabolic}
Ionic gradients are maintained away from electrochemical equilibrium by
continuous metabolic energy input, and the gradients are nonzero.
\end{axiom}

\begin{theorem}[No static equilibrium in an ungrounded excitable graph]
\label{thm:no-equilibrium}
Let $\Gph$ be an ungrounded neuromuscular graph whose vertices carry
voltage-dependent conductances, with total charge conserved and
\Cref{ax:metabolic} in force. Then $\Gph$ admits no static equilibrium: there
is no state with $\dd V^{(v)}/\dd t=0$ at every vertex that is consistent with
the dynamics \eqref{eq:kcl}--\eqref{eq:hh} and with global charge
conservation.
\end{theorem}

\begin{proof}
Suppose such a state exists. Setting $\dd V^{(v)}/\dd t=0$ in \eqref{eq:kcl}
requires, at every vertex, that the net current arriving from neighbours
exactly balance the local ionic current:
$\sum_u I_{uv}=I_{\mathrm{ion}}^{(v)}$ for all $v$.

Consider first a spatially uniform candidate, $V^{(v)}=V^\star$ for all $v$.
Then every inter-vertex current vanishes, so the condition reduces to
$I_{\mathrm{ion}}^{(v)}(V^\star)=0$ at every vertex. By \Cref{ax:metabolic}
the reversal potentials $E_k$ are distinct and the gradients nonzero, so
$I_{\mathrm{ion}}=-\sum_k g_k(V^\star)(V^\star-E_k)$ vanishes only at the
particular potential at which the conductance-weighted driving forces cancel.
That potential is vertex-dependent whenever vertices differ in channel
composition, which is the generic case: motor neurons, muscle fibres, and
receptor endings do not share a conductance profile. Hence no single $V^\star$
satisfies the condition at every vertex, and the uniform candidate fails.

Consider then a non-uniform candidate, $V^{(u)}\ne V^{(v)}$ for some adjacent
$u,v$. The edge between them carries a current proportional to the potential
difference, which is nonzero. Because the graph is ungrounded
(\Cref{prop:no-ground}), this current has no reservoir to enter: by global
charge conservation it must be balanced by redistribution among the remaining
vertices, which changes their charge and hence their potentials, contradicting
the assumption that all time derivatives vanish.

Both cases fail, so no static equilibrium exists.
\end{proof}

\begin{corollary}[Bounded perpetual motion]
\label{cor:perpetual}
Under the hypotheses of \Cref{thm:no-equilibrium}, every trajectory of the
graph is non-static. It is also bounded, since membrane potentials are
confined to the interval spanned by the reversal potentials by the structure
of \eqref{eq:hh}. The trajectory is therefore a bounded non-equilibrium
motion: repeated excursion within a compact region of state space, in the
sense of a dissipative system held from equilibrium by external supply
\citep{nicolis1977,onsager1931}.
\end{corollary}

\begin{remark}[What ``resting'' means]
\label{rem:resting}
\Cref{cor:perpetual} is not a claim that muscles visibly move at rest. It is a
claim about the state of the circuit: the familiar resting potential is a
conditional rest of an isolated cell measured against an external reference,
and in the intact organism it is the centre of a limit cycle rather than a
fixed point. The empirical signature is the tonic discharge observed in
relaxed muscle---spindle afferents at $5$--$30$\,Hz and motor units at
$0.5$--$5$\,Hz \citep{burke1976,matthews1972}---and the continuous
oscillation of the centre of pressure in quiet standing, which never settles
\citep{winter1998,zatsiorsky2000,collins1993}. A language whose execution
model expects a steady state cannot represent either.
\end{remark}

\subsection{Closure is obligatory}

\begin{definition}[Circulation]
\label{def:circulation}
A \emph{circulation} in $\Gph$ is a closed directed path
$v_0\to v_1\to\cdots\to v_n=v_0$ together with a distinguished index $m$
splitting it into an \emph{outbound phase} $v_0\to\cdots\to v_m$ and a
\emph{return phase} $v_m\to\cdots\to v_0$. The outbound phase is efferent and
the return phase afferent.
\end{definition}

\begin{definition}[Aperture]
\label{def:aperture}
A path with a designated outbound phase but no return phase closing it is an
\emph{aperture}. A circuit containing an aperture is \emph{open}; a circuit
all of whose paths close is \emph{closed}.
\end{definition}

We now give the variational statement that selects which closed path a
perturbation follows. Let $J_k$ denote the current on edge $k$ and $g_k$ its
conductance.

\begin{theorem}[Minimum-dissipation closure]
\label{thm:min-dissipation}
In an ungrounded graph with metabolic supply, a charge perturbation introduced
at vertex $u$ redistributes along the closed path $\gamma^\ast$ through $u$
that minimises
\begin{equation}
\Phi[\gamma]\;=\;\sum_{k\in\gamma}\frac{J_k^{2}}{g_k},
\label{eq:dissipation}
\end{equation}
subject to current conservation at every vertex.
\end{theorem}

\begin{proof}
The quantity \eqref{eq:dissipation} is the Joule dissipation associated with
the redistribution, and has units of power. Because the graph is ungrounded,
the perturbation cannot be absorbed and the admissible redistributions are
exactly the closed paths through $u$; this is the constraint set. Minimising a
positive-definite quadratic form over a linear constraint set has a unique
solution characterised by the stationarity conditions, which here state that
the conductance-weighted currents balance at every vertex---that is, that the
solution satisfies Kirchhoff's laws on the closed path. Physically, a
dissipative system under stationary constraint evolves toward the minimum
entropy production consistent with those constraints \citep{onsager1931,%
nicolis1977}, which is the same condition.
\end{proof}

\begin{corollary}[Closure obligation]
\label{cor:closure-obligatory}
A perturbation originating at a central vertex $u$ does not terminate at a
peripheral vertex $w$ unless $w$ lies on a closed path returning to $u$. If no
such path exists, the constraint set of \Cref{thm:min-dissipation} is empty
and no redistribution satisfies the stationarity condition.
\end{corollary}

\begin{proof}
Immediate: the minimisation in \Cref{thm:min-dissipation} is over closed paths
through $u$. If $w$ lies on no such path, no admissible redistribution
terminates at $w$.
\end{proof}

\Cref{cor:closure-obligatory} is the formal content of the design commitment
announced in \Cref{sec:intro}. An efferent projection is not a terminus. It is
one phase of a circulation whose other phase must exist for the perturbation
to resolve. This is why severing a return pathway is not equivalent to
reducing a gain: reducing a gain leaves the constraint set nonempty and
changes the minimiser; severing empties the constraint set.

\begin{principle}[Structural manipulation is not parametric manipulation]
\label{prin:structural}
Modifying a conductance, delay, or gain of a circulation preserves its closure
class and changes the solution of \Cref{thm:min-dissipation}. Removing a
return phase changes the closure class and empties the constraint set. A
language for musculoskeletal experiments must distinguish these
syntactically.
\end{principle}

\subsection{Time-scale strata}

The neuromuscular graph carries loops of widely differing traversal time. We
group them into three strata, whose boundaries are set by measured conduction
and integration times rather than by convention.

\begin{definition}[Strata]
\label{def:strata}
The \emph{reflex} stratum comprises loops closing within $30$--$100$\,ms; the
\emph{spinal} stratum, loops closing within $0.2$--$1$\,s; the
\emph{supraspinal} stratum, loops closing within $1$--$5$\,s.
\end{definition}

The reflex figure is not assumed. The monosynaptic stretch reflex loop
comprises spindle mechanotransduction ($1$--$5$\,ms), afferent conduction
($5$--$15$\,ms at $\sim\!80$\,m/s over limb-dependent distance), monosynaptic
delay ($0.5$--$1$\,ms), efferent conduction ($5$--$15$\,ms), neuromuscular
transmission ($0.5$--$1$\,ms \citep{katz1967,ruff2011}), and
excitation--contraction coupling ($10$--$20$\,ms), summing to $22$--$57$\,ms
for the upper limb and $32$--$72$\,ms for the lower. Measured latencies are
$30$--$40$\,ms and $40$--$50$\,ms respectively \citep{matthews1990},
consistent with the sum.

\begin{proposition}[Spectral separability]
\label{prop:separable}
If the three strata of \Cref{def:strata} are coupled with delays small
relative to their own time constants, the composite trajectory exhibits three
distinguishable spectral bands, with coupling contributing power at sum and
difference frequencies that is small when the time constants are separated by
a factor of five or more.
\end{proposition}

\begin{proof}[Justification]
For weakly coupled oscillators with well-separated natural frequencies, the
composite spectrum is to leading order the sum of the component spectra, with
interaction terms at $\abs{f_i\pm f_j}$ whose amplitude scales with the
coupling strength divided by the frequency separation
\citep{kuramoto1984,strogatz2015}. The strata of \Cref{def:strata} have
central time constants in ratio approximately $1:5:30$, placing them in the
separated regime.
\end{proof}

\Cref{prop:separable} is what makes the stratum discipline of
\Cref{sec:strata} enforceable and its violation observable: a program that
couples strata directly produces cross-band structure that would otherwise be
absent.

% =====================================================================
\section{Lexical structure and grammar}
\label{sec:grammar}
% =====================================================================

\subsection{Token classes}

\begin{definition}[Tokens]
\label{def:tokens}
The lexer recognises:
\emph{keywords}
\texttt{circuit}, \texttt{compartment}, \texttt{capacitance},
\texttt{element}, \texttt{conducts}, \texttt{outbound}, \texttt{return},
\texttt{stratum}, \texttt{delay}, \texttt{gain}, \texttt{floor},
\texttt{derived}, \texttt{event}, \texttt{type},
\texttt{experiment}, \texttt{intact}, \texttt{lesion}, \texttt{observe},
\texttt{without}, \texttt{with}, \texttt{scaling}, \texttt{noise},
\texttt{across}, \texttt{compare}, \texttt{report}, \texttt{let},
\texttt{module}, \texttt{import}, \texttt{reflex}, \texttt{spinal},
\texttt{supraspinal};
\emph{identifiers} matching \texttt{[A-Za-z\_][A-Za-z0-9\_]*} that are not
keywords;
\emph{numeric literals} in decimal or scientific notation, optionally
suffixed by a unit token (\texttt{s}, \texttt{ms}, \texttt{Hz},
\texttt{F}, \texttt{uF}, \texttt{mF}, \texttt{N}, \texttt{m}, \texttt{mm},
\texttt{W}, \texttt{C});
\emph{string literals} delimited by double quotes;
\emph{punctuation} \texttt{\{ \} ( ) [ ] , ; : = -> .};
and \emph{comments} from \texttt{--} to end of line, discarded.
Whitespace is insignificant except as a token separator.
\end{definition}

\subsection{Context-free grammar}

\begin{grammar}
<program> ::= <module-decl>? <import>* <decl>*

<module-decl> ::= `module' <ident> `;'

<import> ::= `import' <ident> `;'

<decl> ::= <compartment-decl> | <circuit-decl>
        \alt <event-type-decl> | <experiment-decl> | <let-decl>

<compartment-decl> ::= `compartment' <ident> `\{'
  `capacitance' `:' <quantity> `;'
  `stratum' `:' <stratum> `;' `\}'

<stratum> ::= `reflex' | `spinal' | `supraspinal'

<circuit-decl> ::= `circuit' <ident> `\{'
  `floor' `:' <floor-spec> `;'
  `outbound' `:' <path> `;'
  `return' `:' <path> `;'
  <element-decl>* `\}'

<floor-spec> ::= `derived' `(' <expr> `)' | <quantity>

<path> ::= <ident> ( `->' <ident> )*

<element-decl> ::= `element' <ident> `conducts' <ident> `->' <ident>
  ( `delay' <quantity> )? ( `gain' <number> )? `;'

<event-type-decl> ::= `event' `type' <ident> `=' <expr> `;'

<experiment-decl> ::= `experiment' <ident> `\{'
  `intact' `:' <circuit-expr> `;'
  <lesion-decl>*
  `observe' `:' <observable-list> `;' `\}'

<lesion-decl> ::= `lesion' <ident> `:' <circuit-expr> `;'

<circuit-expr> ::= <ident>
  \alt <circuit-expr> `without' `return' `(' <ident> `)'
  \alt <circuit-expr> `without' `element' `(' <ident> `)'
  \alt <circuit-expr> `with' <ident> `scaling' <number>
  \alt <circuit-expr> `with' `noise' `across' <stratum> `,' <stratum>
       `(' <number> `)'

<observable-list> ::= <observable> ( `,' <observable> )*

<observable> ::= <ident> ( `(' <arg-list>? `)' )?

<let-decl> ::= `let' <ident> `=' <expr> `;'

<expr> ::= <ident> | <number> | <quantity> | <string>
  \alt <ident> `(' <arg-list>? `)'

<arg-list> ::= <expr> ( `,' <expr> )*

<quantity> ::= <number> <unit>
\end{grammar}

\begin{proposition}[The grammar is LL(1)]
\label{prop:ll1}
The grammar of \Cref{sec:grammar} is LL(1), and therefore unambiguous.
\end{proposition}

\begin{proof}
We verify that one token of lookahead selects a unique production at each
non-terminal. The five \texttt{<decl>} alternatives begin with the pairwise
distinct keywords \texttt{compartment}, \texttt{circuit}, \texttt{event},
\texttt{experiment}, \texttt{let}. Within \texttt{<circuit-decl>} the fields
appear in fixed order and each is introduced by a distinct keyword. The
alternatives of \texttt{<floor-spec>} are distinguished by whether the first
token is the keyword \texttt{derived} or a numeric literal.
\texttt{<circuit-expr>} is left-recursive in the standard sense of a
postfix-operator chain; eliminating the left recursion yields
$\textit{<circuit-expr>}\to\textit{<ident>}\ \textit{<circ-tail>}$ with
$\textit{<circ-tail>}\to$ \texttt{without}\,$\cdots$\,\textit{<circ-tail>}
$\mid$ \texttt{with}\,$\cdots$\,\textit{<circ-tail>} $\mid\ \varepsilon$,
whose alternatives are selected by the distinct keywords \texttt{without} and
\texttt{with} and whose $\varepsilon$-production is taken on any token in the
follow set (\texttt{;} and \texttt{,}), disjoint from both. Within
\texttt{without} the two forms are separated by the distinct keywords
\texttt{return} and \texttt{element}; within \texttt{with}, the \texttt{noise}
form is distinguished from the \texttt{scaling} form by whether the token
following \texttt{with} is the keyword \texttt{noise} or an identifier. The
\texttt{<expr>} alternatives are separated by token class, with the call form
resolved by a single lookahead for \texttt{(}. Hence the grammar is LL(1), and
every LL(1) grammar is unambiguous \citep{aho2006}.
\end{proof}

\begin{remark}[Both phases are mandatory in the grammar]
\label{rem:mandatory-phases}
Note that \texttt{outbound} and \texttt{return} are both non-optional in
\texttt{<circuit-decl>}, whereas \texttt{delay} and \texttt{gain} are
bracketed as optional in \texttt{<element-decl>}. This is the syntactic
expression of \Cref{prin:structural}: a circuit \emph{declaration} always
names both phases, and a circuit lacking a return phase can arise only through
the explicit \texttt{without return} operator inside an experiment, where it
is visible as a deliberate act and receives the diagnostic of
\Cref{sec:aperture}. It is not possible to write an open circuit by
inadvertent omission.
\end{remark}

% =====================================================================
\section{The type system}
\label{sec:typing}
% =====================================================================

\subsection{Types}

\begin{definition}[Types]
\label{def:types}
The types of Vitruvius are
\[
T \;::=\; \Comp \mid \Circ_{c} \mid \Chg_{\gamma} \mid \Obs
\mid \Exp \mid \Res \mid \textsf{Quantity}_{u}
\]
where $c\in\{\textsf{closed},\textsf{open}\}$ indexes the closure class of a
circuit, $\gamma$ ranges over compartment names and indexes a charge by the
compartment whose capacitance it was computed against, and $u$ ranges over
unit expressions.
\end{definition}

Two indices carry the language's non-standard content. The closure index $c$
records whether a circuit is closed or has an aperture, and is computed
statically (\Cref{sec:aperture}). The compartment index $\gamma$ on
$\Chg{}$ records provenance and is the subject of \Cref{sec:compartments}.

A typing environment $\Env$ maps identifiers to types; a stratum environment
$\Delta$ maps circuit elements to strata. The judgement
$\Env;\Delta\vdash e:T\,!\,S$ reads: in environment $\Env$ with stratum
assignment $\Delta$, expression $e$ has type $T$ and touches the set of strata
$S$. The effect component $S$ supports \Cref{sec:strata}.

\subsection{Standard rules}

\begin{align*}
&\textsc{T-Var}\quad
\frac{x:T\in\Env}{\Env;\Delta\vdash x:T\,!\,\emptyset}
\qquad\qquad
\textsc{T-Comp}\quad
\frac{\Env\vdash \mathit{cap}:\textsf{Quantity}_{\mathrm{F}}\quad
      s\in\{\textsf{reflex},\textsf{spinal},\textsf{supraspinal}\}}
     {\Env;\Delta\vdash \texttt{compartment}\ \gamma\ \{\mathit{cap};s\}:\Comp\,!\,\{s\}}
\\[6pt]
&\textsc{T-Elem}\quad
\frac{\Env\vdash u,v:\Comp\quad
      \Env\vdash d:\textsf{Quantity}_{\mathrm{s}}\quad
      \Env\vdash g:\textsf{Quantity}_{1}}
     {\Env;\Delta\vdash \texttt{element}\ k\ \texttt{conducts}\ u\to v\
      \texttt{delay}\ d\ \texttt{gain}\ g:\Obs\,!\,\{\Delta(u),\Delta(v)\}}
\end{align*}

\subsection{Rule I: compartment consistency}
\label{sec:compartments}

Charge is not a scalar in this language. The quantity of charge redistributed
per second across a compartment carrying metabolic power $P$ is obtained from
the energy stored on a capacitance $C$,
\begin{equation}
U=\frac{Q^{2}}{2C}
\quad\Longrightarrow\quad
Q(P;C)=\sqrt{2\,C\,P\,\Delta t},\qquad \Delta t=1\,\mathrm{s},
\label{eq:charge}
\end{equation}
so a charge figure is meaningless without the capacitance against which it was
computed. Aggregate capacitances differ by more than an order of magnitude
across functional compartments: taking specific capacitance
$c_m\approx1\,\mu\mathrm{F/cm}^2$ \citep{cole1941,gentet2000} and the membrane
areas of the relevant populations \citep{kandel2013,purves2018}, a cortical
aggregate is of order $10^{-3}$\,F while a motor-circuit aggregate---motor
neurons, neuromuscular junctions, and sarcolemma of the large postural
muscles---is of order $10^{-4}$\,F.

Adding a charge computed at one capacitance to a charge computed at another is
therefore an error of the same kind as adding a length to a mass.

\begin{definition}[Compartment consistency rule]
\label{def:comp-rule}
\[
\textsc{T-Charge-Add}\quad
\frac{\Env;\Delta\vdash q_1:\Chg_{\gamma}\,!\,S_1\qquad
      \Env;\Delta\vdash q_2:\Chg_{\gamma}\,!\,S_2}
     {\Env;\Delta\vdash q_1+q_2:\Chg_{\gamma}\,!\,S_1\cup S_2}
\]
Charges combine only at a common compartment index. There is no rule
concluding $\Chg_{\gamma_1}+\Chg_{\gamma_2}$ for $\gamma_1\ne\gamma_2$.
\end{definition}

\begin{definition}[Conversion]
\label{def:conversion}
\[
\textsc{T-Charge-Conv}\quad
\frac{\Env;\Delta\vdash q:\Chg_{\gamma}\,!\,S\qquad
      \Env\vdash \gamma,\gamma':\Comp}
     {\Env;\Delta\vdash \texttt{via\_power}(q,\gamma,\gamma'):\Chg_{\gamma'}\,!\,S}
\]
The only route between compartment indices is through power: $q$ is inverted
to a power by \eqref{eq:charge} at $C_{\gamma}$ and re-expressed at
$C_{\gamma'}$. The conversion is explicit in the source.
\end{definition}

\begin{proposition}[No silent mixing]
\label{prop:no-mixing}
If $\Env;\Delta\vdash e:\Chg_{\gamma}\,!\,S$ is derivable, then every charge
literal occurring in $e$ either carries index $\gamma$ or occurs within an
application of \texttt{via\_power} whose target index is $\gamma$.
\end{proposition}

\begin{proof}
By induction on the derivation. The only rules concluding a $\Chg{}$ type are
\textsc{T-Var}, \textsc{T-Charge-Add}, and \textsc{T-Charge-Conv}.
\textsc{T-Var} introduces a literal with its own index. \textsc{T-Charge-Add}
requires both premises to carry the same index as the conclusion, so by the
induction hypothesis every literal in either subterm satisfies the claim.
\textsc{T-Charge-Conv} concludes index $\gamma'$ from a subterm of index
$\gamma$, and the subterm occurs within the application, satisfying the second
disjunct. No other rule concludes $\Chg{}$, so the induction is complete.
\end{proof}

\subsection{Rule II: stratum containment}
\label{sec:strata}

By \Cref{prop:separable} the three strata occupy separated frequency bands.
Physiologically, an element at the reflex stratum does not read supraspinal
state directly: influence from higher levels arrives through the intervening
spinal stratum, with the delays that implies. A program in which a reflex-level
element reads supraspinal state describes a shortcut that the anatomy does not
provide.

We enforce this with an effect discipline. Let $\prec$ order the strata
$\textsf{reflex}\prec\textsf{spinal}\prec\textsf{supraspinal}$, and write
$\mathrm{adj}(s)$ for the strata at distance at most one from $s$ under
$\prec$.

\begin{definition}[Stratum containment rule]
\label{def:stratum-rule}
\[
\textsc{T-Stratum}\quad
\frac{\Env;\Delta\vdash e:T\,!\,S\qquad
      \Delta(k)=s\qquad S\subseteq\mathrm{adj}(s)}
     {\Env;\Delta\vdash \texttt{element}\ k\ \texttt{:=}\ e:\Obs\,!\,\{s\}}
\]
An element assigned to stratum $s$ may touch only strata adjacent to $s$. An
element whose body touches a stratum two or more levels away is rejected.
\end{definition}

\begin{proposition}[Containment is transitive-closure-free]
\label{prop:containment}
Under \textsc{T-Stratum}, influence from stratum $s$ to stratum $s'$ with
$\abs{s-s'}\ge2$ is expressible only as a composition of elements each of
which satisfies the adjacency condition. Consequently every such influence
passes through the intervening stratum and incurs its delay.
\end{proposition}

\begin{proof}
Suppose an element $k$ at stratum $s$ has an effect set containing $s'$ with
$\abs{s-s'}\ge2$. Then $s'\notin\mathrm{adj}(s)$ and the premise
$S\subseteq\mathrm{adj}(s)$ of \textsc{T-Stratum} fails, so no derivation
assigns $k$ a type. Hence any well-typed program realising influence between
$s$ and $s'$ does so through a chain of elements $k_1,\dots,k_n$ with
consecutive strata adjacent, and each $k_i$ contributes its declared delay.
\end{proof}

\begin{remark}[The violation is a model, not merely an error]
\label{rem:violation-is-model}
Direct cross-stratum coupling is precisely what is lost when the structure
that regulates the phase relationship between levels is damaged. The language
therefore provides a construct that introduces such coupling deliberately:
\texttt{with noise across supraspinal, reflex($\eta$)} inserts a
phase-perturbing edge between the named strata with amplitude $\eta$. This
expression is well typed---the escape is explicit and confined to the
experiment section---while an element that quietly reads across two strata is
not. The distinction the language draws is between a modelled lesion and an
unnoticed shortcut.
\end{remark}

\subsection{Rule III: floor positivity}
\label{sec:floor-rule}

Several observables in \Cref{sec:estimation} are expressed as fractions of the
distance between a current value and an irreducible lower bound $\flo$. That
bound must be strictly positive, or the fractions are undefined.

\begin{definition}[Separation cost]
\label{def:sepcost}
Let a circuit be represented as a weighted undirected graph with a
distinguished vertex $\med$ standing for the remainder of the organism. For an
element $v$, the \emph{separation cost} $\sep(v)$ is the minimum total weight
of a set of edges whose removal disconnects $v$ from $\med$; a minimising set
is the \emph{resting cut} of $v$.
\end{definition}

By the max-flow--min-cut correspondence \citep{fordfulkerson1956} this
minimum equals the maximum flow from $v$ to $\med$, and by Menger's theorem
\citep{menger1927} it counts, in the unweighted case, the edge-disjoint paths
that must be severed to isolate $v$. Both readings support the same
interpretation: $\sep(v)$ is the commitment required to hold $v$ apart from
the rest of the system.

\begin{proposition}[Positive floor from connectivity]
\label{prop:floor}
Let a circuit be closed in the sense of \Cref{def:aperture} and connected.
Then $\sep(v)>0$ for every element $v$, and
$\flo:=\min_{v}\sep(v)>0$.
\end{proposition}

\begin{proof}
Since the graph is connected, $v$ is joined to $\med$ by at least one path, so
every disconnecting set is nonempty. Edge weights are positive by
construction---an edge records a conductive pathway of nonzero conductance---so
every disconnecting set has positive total weight, whence $\sep(v)>0$. The
minimum over the finite vertex set of finitely many positive numbers is
positive.
\end{proof}

\begin{definition}[Floor positivity rule]
\label{def:floor-rule}
\[
\textsc{T-Floor}\quad
\frac{\Env;\Delta\vdash C:\Circ_{c}\,!\,S\qquad \flo(C)>0}
     {\Env;\Delta\vdash \texttt{floor}(C):\textsf{Quantity}_{\mathrm{C}}\,!\,S}
\]
A circuit declaring a floor of zero, or one whose derived floor evaluates to
zero, does not type.
\end{definition}

\begin{remark}[Two floor estimators, only one falsifiable]
\label{rem:floor-estimators}
A \texttt{floor} field may be discharged by a literal or by
\texttt{derived(\ldots)}. Two derivations are common and they differ in kind.
The first returns the smallest observed value of the quantity in a finite
sample. This is positive whenever the sample is nonempty, and therefore cannot
return a value inconsistent with a positivity claim: it has no capacity to
falsify. It is moreover biased upward, a sample minimum being an upper bound
on the infimum of the underlying process. The second estimates the asymptote
of separation cost as the model is elaborated, and can return a value
statistically indistinguishable from zero, which would be a falsification.
Vitruvius requires the derivation to be written in the source so that a reader
can see which was used, and \Cref{prop:floor-sensitivity} quantifies what is
at stake.
\end{remark}

\begin{proposition}[Floor sensitivity]
\label{prop:floor-sensitivity}
Let the true floor be $\flo$ and suppose an analysis uses
$\flo'=\flo+\epsilon$. Then a quantity reported as a fraction of the gap,
$\kap=(S_{\mathrm{b}}-S_{\mathrm{a}})/(S_{\mathrm{b}}-\flo)$, is computed as
\[
\kap' \;=\; \kap\cdot\frac{S_{\mathrm{b}}-\flo}{S_{\mathrm{b}}-\flo-\epsilon}
\;=\;\kap\Bigl(1+\frac{\epsilon}{S_{\mathrm{b}}-\flo}\Bigr)+O(\epsilon^{2}),
\]
so the relative error is $\epsilon/(S_{\mathrm{b}}-\flo)$ to first order.
Misestimating the floor inflates the reported values for events acting near
the floor and leaves those acting far from it nearly unchanged.
\end{proposition}

\begin{proof}
Substituting $\flo'$ leaves the numerator unchanged and replaces the
denominator by $S_{\mathrm{b}}-\flo-\epsilon$; the ratio and its first-order
expansion follow.
\end{proof}

The consequence is practical: an upward-biased floor inflates exactly the
quantities that describe the events closest to the irreducible limit, which
are usually the events of interest.

\subsection{Rule IV: typed-event estimation}

The fourth non-standard rule concerns how the effect of an event is estimated,
and requires the development of \Cref{sec:estimation}. We state it here for
completeness and prove its necessity there.

\begin{definition}[Typed estimation rule]
\label{def:typed-rule}
\[
\textsc{T-Event}\quad
\frac{\Env;\Delta\vdash \gamma:\Obs\,!\,S\qquad
      \Env\vdash \theta:\textsf{EventType}\qquad
      \abs{\mathcal{T}_{\theta}}\ge2}
     {\Env;\Delta\vdash \kap(\gamma;\theta):\textsf{Quantity}_{1}\,!\,S}
\]
The catalytic power of an event is well typed only when the event carries a
declared type $\theta$ possessing at least two recorded instances. There is no
rule assigning a power to an untyped event.
\end{definition}

% =====================================================================
\section{The aperture analysis}
\label{sec:aperture}
% =====================================================================

We now specify the static analysis that gives the language its purpose: the
decision, before any integration, of whether a program's circuits are closed,
and the diagnostic issued when they are not.

\begin{definition}[Closure index]
\label{def:closure-index}
For a circuit expression $e$, the \emph{closure index} $c(e)$ is
$\textsf{closed}$ if every declared outbound path in $e$ is matched by a
return path from its terminal compartment to its initial compartment, and
$\textsf{open}$ otherwise.
\end{definition}

\begin{theorem}[Closure is decidable in linear time]
\label{thm:closure-decidable}
The closure index of a circuit expression is decidable, and computable in time
linear in the number of elements.
\end{theorem}

\begin{proof}
A circuit declaration names its outbound path $p$ and return path $r$
explicitly (\Cref{sec:grammar}), so closure of a base circuit is checked by
comparing the terminal compartment of $p$ with the initial compartment of $r$
and the initial of $p$ with the terminal of $r$: two comparisons.

For a compound expression we proceed by structural induction. The operators
\texttt{with $k$ scaling $\lambda$} and
\texttt{with noise across $s_1,s_2(\eta)$} alter weights and add a coupling
edge respectively, and neither removes an element from a path, so
$c(e\ \texttt{with}\ \cdots)=c(e)$. The operator
\texttt{without element($k$)} removes $k$; the result is open if $k$ lies on
the return path of any circulation of $e$ and closed otherwise, decided by a
single membership test. The operator \texttt{without return($a$)} removes the
return phase named $a$ and yields $\textsf{open}$ by definition.

Each case performs work bounded by a constant plus one traversal of the
element list, so the total is linear in the number of elements.
\end{proof}

\begin{definition}[Aperture report]
\label{def:aperture-report}
When $c(e)=\textsf{open}$, the compiler emits an \emph{aperture report} naming
(i) the circulation left open, (ii) the outbound phase now without a return,
and (iii) the prediction of \Cref{cor:closure-obligatory}: that the
perturbation admits no closed redistribution, so the expected behaviour is
failure of the movement to resolve, not increased variance about a preserved
trajectory.
\end{definition}

\begin{remark}[The report is not an error]
\label{rem:report-not-error}
An open circuit is a legitimate object of study---it is the model of
deafferentation, and simulating it is the point of the exercise. The aperture
report is therefore a diagnostic, not a rejection: the program compiles, runs,
and returns results. What the language refuses is silence. A user who opens a
circulation is told so, in terms that state what the framework predicts, before
the numerics begin. This is the discipline the language exists to impose, and
the reason \Cref{prin:structural} is enforced syntactically rather than left to
a comment.
\end{remark}

\begin{proposition}[Attenuation preserves closure]
\label{prop:attenuation}
For any circuit $e$, element $k$, and scaling $\lambda>0$,
$c(e\ \texttt{with}\ k\ \texttt{scaling}\ \lambda)=c(e)$. In particular no
finite attenuation of a return element opens a circulation.
\end{proposition}

\begin{proof}
By \Cref{thm:closure-decidable}, the scaling operator alters edge weights and
removes no element, so the outbound and return paths are unchanged as vertex
sequences and the matching test returns the same answer. That $\lambda>0$ is
required is what makes the statement non-vacuous: the grammar admits no
$\lambda=0$, and the elimination of an element is available only through
\texttt{without}, which is the syntactically distinct structural operator.
\end{proof}

\Cref{prop:attenuation} is the formal content of the claim that Vitruvius
separates structural from parametric manipulation. In a signal-and-gain
architecture, $\lambda\to0$ and severance are the same limit; here they are
different operators with different closure indices, and the difference is
visible statically.

% =====================================================================
\section{Operational semantics}
\label{sec:semantics}
% =====================================================================

\subsection{Configurations}

\begin{definition}[Configuration]
\label{def:config}
A \emph{configuration} is a quadruple
$\langle e,\ K,\ \store,\ \rho\rangle$ in which $e$ is the term under
evaluation, $K$ is the circuit under study, $\store$ is the observation store
mapping observable names to recorded values, and $\rho\in\N$ is the
\emph{committed record}, the count of observations discharged.
\end{definition}

\subsection{Reduction rules}

We write $\mathcal{B}(K,o)\Downarrow x$ for ``the backend, given circuit $K$,
discharges observable $o$ to value $x$'' (\Cref{sec:backends}). Let $O$ denote
the pending observable set of an experiment and $L$ its lesion list.

\begin{align*}
&\textsc{E-Lesion}\quad
\frac{\ell\in L\qquad K'=\mathrm{apply}(\ell,K)}
{\bigl\langle \experiment\{K;L;O\},\store,\rho\bigr\rangle
\steps
\bigl\langle \experiment\{K';L\setminus\{\ell\};O\},\store,\rho\bigr\rangle}
\\[8pt]
&\textsc{E-Observe}\quad
\frac{L=\emptyset\qquad o\in O\qquad \mathcal{B}(K,o)\Downarrow x}
{\bigl\langle \experiment\{K;\emptyset;O\},\store,\rho\bigr\rangle
\steps
\bigl\langle \experiment\{K;\emptyset;O\setminus\{o\}\},
 \store[o\mapsto x],\rho+1\bigr\rangle}
\\[8pt]
&\textsc{E-Report}\quad
\frac{L=\emptyset\qquad O=\emptyset}
{\bigl\langle \experiment\{K;\emptyset;\emptyset\},\store,\rho\bigr\rangle
\steps
\bigl\langle \Res(\store),\store,\rho\bigr\rangle}
\end{align*}

\textsc{E-Lesion} applies one pending ablation; \textsc{E-Observe} discharges
one pending observable and increments the record; \textsc{E-Report} produces
the result once both sets are exhausted. Lesions precede observations, so an
experiment observes the circuit it declared and not an intermediate.

\subsection{Metatheory}

\begin{theorem}[Progress]
\label{thm:progress}
If $\Env;\Delta\vdash e:\Exp\,!\,S$ and $e$ is not of the form $\Res(\cdot)$,
then for any store $\store$ and record $\rho$ there is a configuration to
which $\langle e,K,\store,\rho\rangle$ steps.
\end{theorem}

\begin{proof}
A well-typed non-value term of type $\Exp$ is an experiment, that being the
only form for which a rule concludes $\Exp$. Consider its lesion list $L$ and
observable set $O$. If $L\ne\emptyset$, pick $\ell\in L$; the operator
$\mathrm{apply}$ is total on well-typed lesions, since each of the four
circuit operators is defined for every circuit (removal of an absent element
being the identity), so \textsc{E-Lesion} applies. If $L=\emptyset$ and
$O\ne\emptyset$, pick $o\in O$; the backend obligation
(\Cref{def:backend}) requires $\mathcal{B}$ to be total on well-typed
observables, so \textsc{E-Observe} applies. If both are empty,
\textsc{E-Report} applies. The three cases are exhaustive, so the term is not
stuck.
\end{proof}

\begin{theorem}[Preservation]
\label{thm:preservation}
If $\Env;\Delta\vdash e:T\,!\,S$ and
$\langle e,K,\store,\rho\rangle\steps\langle e',K',\store',\rho'\rangle$, then
$\Env;\Delta\vdash e':T\,!\,S'$ for some $S'\subseteq S$.
\end{theorem}

\begin{proof}
By cases on the rule applied. \textsc{E-Lesion} yields an experiment differing
only in its circuit and a smaller lesion list. The typing of an experiment
requires its circuit to be well typed at $\Circ_{c}$ for some $c$, and each
circuit operator maps well-typed circuits to well-typed circuits---possibly
altering the closure index, which is a change of index rather than of type
constructor and is admitted by \Cref{def:types}. The effect set does not grow,
since no operator introduces an element outside the strata already present
except the \texttt{noise} operator, whose strata are named in its own typing
premise and hence already in $S$. \textsc{E-Observe} yields an experiment with
a smaller observable set and an extended store, neither of which appears in
the typing premises for $\Exp$. \textsc{E-Report} yields $\Res(\store)$, of
type $\Res$; the experiment's declared result type is $\Res$ by
\textsc{T-Exp}, so the type is preserved. In every case the type constructor is
unchanged and the effect set does not grow.
\end{proof}

\begin{corollary}[Soundness]
\label{cor:soundness}
A well-typed program does not get stuck: it either reduces indefinitely or
reaches a value of type $\Res$. With \Cref{thm:termination} it reaches such a
value in finitely many steps.
\end{corollary}

\begin{proof}
Immediate from \Cref{thm:progress,thm:preservation} by the standard syntactic
argument \citep{wright1994}.
\end{proof}

\begin{theorem}[Lesion composition is a join-semilattice operation]
\label{thm:lesion-monoid}
The application of lesions is associative, commutative, and idempotent on
circuits: for lesions $\ell_1,\ell_2$ and circuit $K$,
\begin{align*}
\mathrm{apply}(\ell_1,\mathrm{apply}(\ell_2,K))
&=\mathrm{apply}(\ell_2,\mathrm{apply}(\ell_1,K)),\\
\mathrm{apply}(\ell,\mathrm{apply}(\ell,K))&=\mathrm{apply}(\ell,K).
\end{align*}
Consequently the lesioned circuit of an experiment is independent of the order
in which its lesions are written.
\end{theorem}

\begin{proof}
We treat the operators in turn. \texttt{without element($k$)} and
\texttt{without return($a$)} both delete a named item from the element set;
deletion of distinct names commutes, deletion of the same name twice equals
deletion once, and deletion is associative because set difference is. So both
removal operators are commutative, idempotent, and associative among
themselves and with each other.

\texttt{with $k$ scaling $\lambda$} multiplies the weight of $k$ by $\lambda$.
Two scalings of distinct elements commute since they touch disjoint weights.
Two scalings of the same element compose multiplicatively, and are therefore
commutative but idempotent only for $\lambda=1$. We accordingly define the
lesion list to be a \emph{set} of scalings keyed by element name, with
duplicate keys rejected at typecheck; under that restriction at most one
scaling applies per element and idempotence holds. Scaling commutes with
removal in the following sense: if $k$ is removed then its weight is
irrelevant, so scaling before or after removal yields the same circuit; if $k$
is not removed, the operations touch different items.

\texttt{with noise across $s_1,s_2(\eta)$} adjoins a coupling edge. Adjunction
of edges commutes and is idempotent when edges are identified by their
endpoint pair and amplitude, and duplicate declarations are rejected as above.
Adjunction commutes with scaling (disjoint items) and with removal (a coupling
edge is not on a declared path, so its presence does not affect which elements
are deleted).

All pairs commute and each operator is idempotent under the stated key
restriction, so the induced operation on circuits is a commutative idempotent
monoid action; the resulting circuit is therefore the join of the individual
lesions and independent of order by induction on the length of the list.
\end{proof}

\begin{theorem}[Termination]
\label{thm:termination}
Every well-typed experiment over a finite lesion list and finite observable
set reaches a $\Res$ value in at most $\abs{L}+\abs{O}+1$ steps.
\end{theorem}

\begin{proof}
\textsc{E-Lesion} strictly decreases $\abs{L}$ and leaves $\abs{O}$ unchanged;
\textsc{E-Observe} applies only when $L=\emptyset$ and strictly decreases
$\abs{O}$; \textsc{E-Report} applies once. The lexicographic measure
$(\abs{L},\abs{O})$ therefore strictly decreases at every step of the first
two kinds and is bounded below, so after at most $\abs{L}$ lesion steps and
$\abs{O}$ observation steps both are zero and \textsc{E-Report} fires, giving
the stated bound.
\end{proof}

\begin{theorem}[Monotone record and non-return]
\label{thm:monotone}
The record $\rho$ is non-decreasing along any reduction and strictly increases
at each \textsc{E-Observe}. Two configurations with identical terms, circuits,
and stores but different records are distinct; hence no experiment returns to
a previous configuration.
\end{theorem}

\begin{proof}
\textsc{E-Observe} increments $\rho$; the other rules leave it unchanged, and
no rule decrements it. Configurations with records $r<r'$ are separated by the
predicate $\rho\le r$, hence distinct. A return would require $\rho$ to
decrease, which no rule permits.
\end{proof}

\begin{remark}[Why the record matters here]
\label{rem:record}
By \Cref{cor:perpetual} the underlying system has no terminal state, so a run
cannot be individuated by its final configuration in the way a conventional
program can. The record supplies an intrinsic ordering of observations that
does not appeal to a wall clock and does not presume convergence, and it is
what makes a repeated value in the store distinguishable from a return to an
earlier state.
\end{remark}

% =====================================================================
\section{Estimating the effect of an event}
\label{sec:estimation}
% =====================================================================

Experiments in this language routinely ask how much a given event contributes
to a change. This section shows that the natural way to ask is vacuous, and
specifies the way that is not. The result governs rule IV of
\Cref{sec:typing}.

\subsection{Catalytic power}

\begin{definition}[Catalytic power]
\label{def:power}
Let an event $\gamma$ carry the system from a state with observable value
$S_{\mathrm{b}}>\flo$ to one with value $S_{\mathrm{a}}$, where $\flo$ is the
floor of \Cref{sec:floor-rule}. Its \emph{catalytic power} is
\begin{equation}
\kap(\gamma)\;=\;\frac{S_{\mathrm{b}}-S_{\mathrm{a}}}{S_{\mathrm{b}}-\flo},
\label{eq:kappa}
\end{equation}
the fraction of the outstanding distance to the floor that the event closes.
Its \emph{residual factor} is
$\bar\kap(\gamma)=1-\kap(\gamma)
=(S_{\mathrm{a}}-\flo)/(S_{\mathrm{b}}-\flo)$.
\end{definition}

\begin{theorem}[Multiplicative composition]
\label{thm:composition}
For a sequence of events $\gamma_1,\dots,\gamma_n$ applied in order over a
fixed floor, with the post-state of each the pre-state of the next, the net
power satisfies
\begin{equation}
\kap_{\mathrm{net}}
\;=\;1-\prod_{i=1}^{n}\bigl(1-\kap(\gamma_i)\bigr).
\label{eq:composition}
\end{equation}
\end{theorem}

\begin{proof}
Write $g_i=S_i-\flo$ for the gap before event $i$. By
\Cref{def:power}, $\bar\kap(\gamma_i)=g_{i+1}/g_i$. Hence
\[
\prod_{i=1}^{n}\bar\kap(\gamma_i)=\prod_{i=1}^{n}\frac{g_{i+1}}{g_i}
=\frac{g_{n+1}}{g_1}=\frac{S_{n+1}-\flo}{S_1-\flo}=1-\kap_{\mathrm{net}},
\]
which rearranges to \eqref{eq:composition}.
\end{proof}

\subsection{The degenerate estimator}

\Cref{thm:composition} suggests an immediate empirical test: estimate each
event's power, form the product, compare with the directly measured net
change. We show this test cannot fail.

\begin{definition}[Instance-specific estimation]
\label{def:instance}
Powers are estimated \emph{instance-specifically} when each $\kap(\gamma_i)$
is computed from \eqref{eq:kappa} using the same observed values
$S_i,S_{i+1}$ that determine the measured net change of the sequence they
belong to.
\end{definition}

\begin{theorem}[The instance-specific test is an identity]
\label{thm:telescoping}
Under instance-specific estimation, the composed prediction and the measured
net power are equal for every sequence, identically:
\[
1-\prod_{i=1}^{n}\bigl(1-\kap(\gamma_i)\bigr)
\;=\;\frac{S_1-S_{n+1}}{S_1-\flo}
\;=\;\kap_{\mathrm{measured}}.
\]
Consequently any discrepancy statistic---a correlation, a root-mean-square
error, a regression slope---attains its degenerate value on every data set
whatsoever, independently of the process generating the data.
\end{theorem}

\begin{proof}
This is \Cref{thm:composition} read as an identity rather than as a
prediction. With $g_i=S_i-\flo$, instance-specific estimation gives
$1-\kap(\gamma_i)=g_{i+1}/g_i$ exactly, so the product telescopes to
$g_{n+1}/g_1$, whence
$1-\prod_i(1-\kap_i)=(S_1-S_{n+1})/(S_1-\flo)$, which is the measured
quantity. The two are the same function of the data, so their difference
vanishes pointwise.
\end{proof}

\begin{corollary}[No null hypothesis]
\label{cor:no-null}
There is no data set and no generating process under which the
instance-specific test fails. It therefore has no power against any
alternative, and a reported agreement carries no evidential weight.
\end{corollary}

\begin{remark}[Why this is easy to miss]
\label{rem:easy-to-miss}
The obstruction is invisible in the implementation. The two quantities are
computed by different routines---one iterating over events and multiplying,
the other differencing the endpoints---and they agree to machine precision
across arbitrarily many trials. The agreement presents as a striking
confirmation. It is a restatement of \Cref{thm:composition}. A correlation of
exactly unity between a composed prediction and its own measurement is
therefore an occasion for suspicion rather than satisfaction, and any
framework whose composition law telescopes will exhibit it.
\end{remark}

\subsection{The typed estimator}

The identity arises because the constituent powers and the aggregate outcome
are functions of the same numbers. Breaking it requires estimating powers from
a source independent of the particular sequence being predicted. The natural
such source is the event's \emph{type}: for a stretch reflex, the ordered pair
of compartments it connects together with its stratum.

\begin{definition}[Type-averaged power]
\label{def:type-averaged}
Let each event carry a type $t(\gamma)$ from a finite label set. For a corpus
of observed events, the type-averaged power of type $\theta$ is
\[
\bar\kap_{\theta}
=\frac{1}{\abs{\mathcal{T}_{\theta}}}
 \sum_{\gamma\in\mathcal{T}_{\theta}}\kap(\gamma),
\qquad
\mathcal{T}_{\theta}=\{\gamma:t(\gamma)=\theta\},
\]
and the type-averaged prediction for a sequence is
$\hat\kap=1-\prod_{i}\bigl(1-\bar\kap_{t(\gamma_i)}\bigr)$.
\end{definition}

\begin{theorem}[The typed estimator is non-degenerate]
\label{thm:nondegenerate}
Suppose some type $\theta$ has at least two instances with distinct powers.
Then there exist sequences for which
$\hat\kap\ne\kap_{\mathrm{measured}}$, so the discrepancy is not identically
zero and the test has a genuine null hypothesis.
\end{theorem}

\begin{proof}
Let $\theta$ have instances with powers $\kap\ne\kap'$; then
$\bar\kap_{\theta}$ differs from at least one of them, say from $\kap$.
Consider the length-one sequence consisting of that instance. Then
$\kap_{\mathrm{measured}}=\kap$ while $\hat\kap=\bar\kap_{\theta}\ne\kap$, so
the discrepancy is nonzero at that point and hence not identically zero as a
function of the data.
\end{proof}

This theorem is the justification for rule IV (\Cref{def:typed-rule}): the
premise $\abs{\mathcal{T}_{\theta}}\ge2$ is exactly the hypothesis of
\Cref{thm:nondegenerate}. A language that permitted untyped power estimation
would permit its users to compute a statistic that cannot fail.

\begin{proposition}[Exact bias of the typed prediction]
\label{prop:bias}
Write $\kap_i=\bar\kap_{t(i)}+\delta_i$. Then
\[
\frac{1-\kap_{\mathrm{measured}}}{1-\hat\kap}
=\prod_{i=1}^{n}\Bigl(1-\frac{\delta_i}{1-\bar\kap_{t(i)}}\Bigr),
\]
so the prediction is exact if and only if the within-type deviations vanish,
and to first order
$\hat\kap-\kap_{\mathrm{measured}}\approx
(1-\hat\kap)\sum_i \delta_i/(1-\bar\kap_{t(i)})$.
\end{proposition}

\begin{proof}
By \Cref{thm:composition},
$1-\kap_{\mathrm{measured}}=\prod_i(1-\bar\kap_{t(i)}-\delta_i)
=\prod_i(1-\bar\kap_{t(i)})\prod_i\bigl(1-\delta_i/(1-\bar\kap_{t(i)})\bigr)$,
and $\prod_i(1-\bar\kap_{t(i)})=1-\hat\kap$ by \Cref{def:type-averaged}.
Dividing gives the ratio; expanding to first order gives the approximation.
\end{proof}

\subsection{When the typed estimator has power}

\begin{definition}[Type separation]
\label{def:eta}
Let $\mathrm{Var}_{\mathrm{w}}$ be the mean within-type variance of $\kap$ and
$\mathrm{Var}_{\mathrm{b}}$ the variance across types of the type means. The
\emph{type separation} is
\[
\eta=\frac{\mathrm{Var}_{\mathrm{b}}}
{\mathrm{Var}_{\mathrm{b}}+\mathrm{Var}_{\mathrm{w}}}\in[0,1].
\]
\end{definition}

\begin{proposition}[Degeneracy at vanishing separation]
\label{prop:degeneracy}
If $\eta=0$, the type-averaged prediction is constant across sequences of a
given length, its sample variance vanishes, its correlation with the measured
outcome is undefined, and the multiplicative law is indistinguishable from any
competitor that is also a function of the type means alone.
\end{proposition}

\begin{proof}
If all $\bar\kap_{\theta}$ equal a common $\bar\kap$, then
$\hat\kap=1-(1-\bar\kap)^{n}$ depends only on $n$. Among sequences of fixed
length the predictor is constant, so its sample variance is zero and the
Pearson correlation, having a zero denominator, is undefined. Any alternative
law computed from the same type means is likewise constant at fixed length and
therefore not separable from it.
\end{proof}

\begin{remark}[Reporting obligation]
\label{rem:report-eta}
\Cref{prop:degeneracy} identifies a failure mode that must be excluded before
a negative result is interpreted. If the chosen typing maps every event into a
narrow band of powers, the test is underpowered by construction regardless of
whether the composition law holds, and a negative result is then evidence
about the typing rather than about the process. Vitruvius therefore reports
$\eta$ alongside any composed estimate, and flags results computed at low
$\eta$. The strata of \Cref{def:strata} are separated by factors of five to
thirty, which is the regime in which $\eta$ should be large; a low $\eta$ under
that typing is itself an informative finding about the model.
\end{remark}

\begin{proposition}[Competing composition laws are ordered]
\label{prop:competitors}
For $n\ge2$ and powers in $(0,1)$,
\[
\max_i\kap_i \;\le\; 1-\prod_i(1-\kap_i)
\;\le\;\min\Bigl(\sum_i\kap_i,\,1\Bigr),
\]
with both inequalities strict unless some $\kap_i=0$; and the geometric-mean
law $1-\bigl(\prod_i(1-\kap_i)\bigr)^{1/n}$ lies strictly below the
multiplicative law whenever $\prod_i(1-\kap_i)<1$.
\end{proposition}

\begin{proof}
For the left inequality, every factor of $\prod_i(1-\kap_i)$ is at most $1$
and one factor equals $1-\max_i\kap_i$, so
$\prod_i(1-\kap_i)\le 1-\max_i\kap_i$; negating gives the claim, with equality
exactly when all other powers vanish. For the right, the Weierstrass
inequality gives $\prod_i(1-\kap_i)\ge 1-\sum_i\kap_i$, whence
$1-\prod_i(1-\kap_i)\le\sum_i\kap_i$; combining with
$\kap_{\mathrm{net}}\le1$ gives the bound. For the geometric mean, with
$P=\prod_i(1-\kap_i)<1$ we have $P^{1/n}>P$ for $n\ge2$, so
$1-P^{1/n}<1-P$.
\end{proof}

\Cref{prop:competitors} supplies the comparison set: the laws are distinct and
ordered, so a corpus with adequate $\eta$ can in principle select among them,
and the multiplicative law being intermediate means a finding in its favour is
two-sided.

% =====================================================================
\section{Backends}
\label{sec:backends}
% =====================================================================

The language computes nothing. Every observable is discharged by a backend,
which is obliged to satisfy the following interface and is otherwise opaque.

\begin{definition}[Backend]
\label{def:backend}
A \emph{backend} $\mathcal{B}$ is a total map from a well-typed circuit and a
well-typed observable to a value, discharging four obligations:
\begin{enumerate}[label=\textup{(B\arabic*)},leftmargin=3.4em,nosep]
\item \textbf{Totality.} $\mathcal{B}(K,o)$ is defined for every well-typed
$K$ and $o$, including open $K$. An open circuit yields a value describing the
failure---a divergence time, an unbounded excursion---not an exception.
\item \textbf{Floor agreement.} $\mathcal{B}$ reports the floor it used, and
it must equal the floor of the circuit declaration up to a stated tolerance.
\item \textbf{Stratum honesty.} $\mathcal{B}$ reports the bandwidth over which
each observable was computed, so that a value attributed to one stratum can be
checked against that stratum's band.
\item \textbf{Determinism modulo seed.} Given a circuit, an observable, and a
seed, $\mathcal{B}$ returns the same value.
\end{enumerate}
\end{definition}

Obligation (B1) deserves emphasis. A backend that raises on a diverging
integration would convert the central experimental result of the language---the
failure of an open circuit to resolve---into a crash, discarding the
measurement that the experiment was constructed to obtain. The divergence is
the finding.

\begin{proposition}[Backend independence of the static analyses]
\label{prop:backend-independent}
The closure index, the compartment analysis, the stratum analysis, and the
floor positivity check are computed without reference to $\mathcal{B}$.
Consequently a program's diagnostics are identical under every conforming
backend.
\end{proposition}

\begin{proof}
By inspection of the rules. \Cref{thm:closure-decidable} computes $c(e)$ from
the element list alone. \textsc{T-Charge-Add} and \textsc{T-Charge-Conv}
inspect compartment indices, which are declared. \textsc{T-Stratum} inspects
$\Delta$, which is declared. \textsc{T-Floor} inspects the floor field, whose
\texttt{derived} form is evaluated over the declared graph by
\Cref{def:sepcost}. None of these consults $\mathcal{B}$, which appears only in
\textsc{E-Observe}.
\end{proof}

\Cref{prop:backend-independent} is what permits a workflow in which a model is
checked cheaply and integrated expensively: the diagnostics that matter for
model correctness are available before any numerical work, and do not change
when the backend does.

\begin{remark}[Three natural backends]
\label{rem:three-backends}
The interface admits backends of quite different character. A
\emph{continuous} backend integrates \eqref{eq:hh} together with cross-bridge
kinetics \citep{huxley1957,zahalak1981} and a musculotendon model
\citep{zajac1989,millard2013}, using standard stiff solvers
\citep{hairer1996,shampine1997}; it is the most expensive and the only one
that returns waveforms. A \emph{discrete} backend represents each compartment
by a finite state set and each element by a transition rule, and returns
occupancies and dwell statistics; it is cheap and adequate for questions about
reachability and closure failure. A \emph{spectral} backend linearises about
the operating cycle and returns band powers directly by periodogram methods
\citep{welch1967}, which suffices for the observables of
\Cref{prop:separable}. A model may be developed against the discrete backend
and only later integrated, since by \Cref{prop:backend-independent} the
diagnostics do not change.
\end{remark}

% =====================================================================
\section{Worked programs}
\label{sec:programs}
% =====================================================================

We exercise every construct on three experiments of increasing scope.

\subsection{A monosynaptic stretch reflex}

The first program declares compartments with their capacitances and strata,
assembles a circulation from spindle to motor neuron to muscle and back, and
observes the loop latency predicted in \Cref{sec:foundations}.

\begin{lstlisting}[style=vvs,caption={\texttt{soleus\_reflex.vvs} --- a closed
monosynaptic loop.},label={lst:reflex}]
module soleus_reflex;

-- Compartments carry a capacitance and a stratum. Capacitance is
-- required: it indexes every charge computed against this compartment.
compartment spindle    { capacitance: 2.0e-9 F;  stratum: reflex; }
compartment ia_afferent{ capacitance: 8.0e-9 F;  stratum: reflex; }
compartment alpha_mn   { capacitance: 1.2e-8 F;  stratum: reflex; }
compartment nmj        { capacitance: 4.0e-9 F;  stratum: reflex; }
compartment fibre      { capacitance: 1.41e-4 F; stratum: reflex; }

circuit soleus_loop {
  -- The floor is derived from the resting cut, not from a sample
  -- minimum; see the discussion of estimator falsifiability.
  floor    : derived(resting_cut(alpha_mn));

  outbound : alpha_mn -> nmj -> fibre;
  return   : fibre -> spindle -> ia_afferent -> alpha_mn;

  element mn_axon   conducts alpha_mn    -> nmj        delay 5.0 ms;
  element endplate  conducts nmj         -> fibre      delay 0.8 ms;
  element ec_couple conducts fibre       -> spindle    delay 15.0 ms;
  element mechano   conducts spindle     -> ia_afferent delay 3.0 ms;
  element ia_axon   conducts ia_afferent -> alpha_mn   delay 8.0 ms
                                                       gain 1.0;
}

experiment latency_budget {
  intact  : soleus_loop;
  observe : loop_latency, tonic_rate, resting_cut_weight;
}
\end{lstlisting}

The declared delays sum to $31.8$\,ms, within the $30$--$40$\,ms band measured
for upper-limb reflexes and just below the $40$--$50$\,ms band for the lower
limb \citep{matthews1990}; the residual is the limb-length-dependent
conduction term, which a backend supplies from the declared segment geometry.
The program typechecks: every element lies at the reflex stratum, so
\textsc{T-Stratum} is satisfied trivially; the circuit is closed, since the
return path begins at \texttt{fibre}, the terminus of the outbound path, and
ends at \texttt{alpha\_mn}, its origin.

\subsection{Ablation of the return phase}

The second program is the language's characteristic experiment. It compares an
intact postural loop against three manipulations: attenuation of the return
element, removal of the return element, and removal of the entire return
phase.

\begin{lstlisting}[style=vvs,caption={\texttt{postural\_ablation.vvs} ---
attenuation and severance are distinct operators.},
label={lst:ablation}]
module postural_ablation;
import soleus_reflex;

compartment spinal_in  { capacitance: 3.0e-8 F; stratum: spinal; }
compartment cortex     { capacitance: 1.0e-3 F; stratum: supraspinal; }

circuit postural_loop {
  floor    : derived(resting_cut(spinal_in));

  outbound : cortex -> spinal_in -> alpha_mn -> nmj -> fibre;
  return   : fibre -> spindle -> ia_afferent -> spinal_in -> cortex;

  element descend   conducts cortex      -> spinal_in   delay 12.0 ms;
  element premotor  conducts spinal_in   -> alpha_mn    delay 2.0 ms;
  element mn_axon   conducts alpha_mn    -> nmj         delay 5.0 ms;
  element endplate  conducts nmj         -> fibre       delay 0.8 ms;
  element ec_couple conducts fibre       -> spindle     delay 15.0 ms;
  element mechano   conducts spindle     -> ia_afferent delay 3.0 ms;
  element ia_axon   conducts ia_afferent -> spinal_in   delay 8.0 ms
                                                        gain 1.0;
  element ascend    conducts spinal_in   -> cortex      delay 14.0 ms;
}

experiment deafferentation {
  intact  : postural_loop;

  -- Parametric: the loop stays closed, the gain falls to one tenth.
  lesion attenuated : postural_loop with ia_axon scaling 0.1;

  -- Structural: the return element is deleted. The compiler reports
  -- an aperture, because no path now returns fibre to spinal_in.
  lesion severed    : postural_loop without element(ia_axon);

  -- Structural: the whole return phase is deleted.
  lesion abolished  : postural_loop without return(fibre);

  observe : cop_rms, divergence_time, band_power(reflex),
            band_power(supraspinal), coupling_index;
}
\end{lstlisting}

The three lesions are not three points on one continuum. By
\Cref{prop:attenuation} the first preserves the closure index and the compiler
is silent; by \Cref{thm:closure-decidable} the second and third set it to
$\textsf{open}$ and each receives an aperture report naming the circulation
and stating the prediction of \Cref{cor:closure-obligatory}. The observable
\texttt{divergence\_time} is well defined only for the open cases; obligation
(B1) requires the backend to return it rather than raise.

This is the point at which the language earns its design. A signal-and-gain
formulation would express all three as gains of $0.1$, $0$, and $0$, making
the second and third identical and the first a small perturbation of them.
Vitruvius separates them statically, and reports which of them is predicted to
fail structurally rather than merely to perform worse.

\subsection{A cross-stratum coupling lesion}

The third program exercises the stratum discipline and the estimation
machinery. It models loss of phase regulation between the supraspinal and
reflex strata---a coupling pathology rather than a conduction failure---and
reports the type separation alongside the composed estimate.

\begin{lstlisting}[style=vvs,caption={\texttt{coupling\_lesion.vvs} --- a
deliberate stratum violation, and the reporting obligation.},
label={lst:coupling}]
module coupling_lesion;
import postural_ablation;

-- Event types are declared so that catalytic power can be estimated
-- from a type rather than from the instance being predicted. Rule IV
-- rejects an untyped power, and each type needs >= 2 instances.
event type reflex_arc     = pair(spindle, alpha_mn);
event type spinal_relay   = pair(spinal_in, alpha_mn);
event type descending_set = pair(cortex, spinal_in);

experiment phase_regulation {
  intact  : postural_loop;

  -- Well typed: the escape is explicit and confined to the experiment.
  -- An *element* reading across two strata would be rejected by
  -- T-Stratum; this construct declares the violation as the model.
  lesion decoupled : postural_loop
                       with noise across supraspinal, reflex (0.35);

  observe : coupling_index,
            band_power(reflex), band_power(spinal),
            band_power(supraspinal),
            kappa(reflex_arc), kappa(spinal_relay),
            type_separation;
}
\end{lstlisting}

Two features of this program are enforced rather than conventional. First, the
cross-stratum edge could not have been written as an element: an element at
the reflex stratum whose body touched \texttt{supraspinal} would violate
\textsc{T-Stratum}, since the two are not adjacent. The only route to such
coupling is the explicit \texttt{with noise across} operator, so a
cross-stratum influence in a Vitruvius model is always a declared lesion and
never an oversight (\Cref{rem:violation-is-model}).

Second, each \texttt{kappa} observable names an event type. Rule IV
(\Cref{def:typed-rule}) rejects a bare \texttt{kappa(mechano)} on a single
instance, which is exactly the instance-specific estimator that
\Cref{thm:telescoping} shows cannot fail. The accompanying
\texttt{type\_separation} observable returns $\eta$ of \Cref{def:eta}, without
which, by \Cref{prop:degeneracy}, a reported agreement or disagreement cannot
be interpreted.

% =====================================================================
\section{Implementation and empirical findings}
\label{sec:implementation}
% =====================================================================

The specification above is a set of claims about what a language can decide
and what it must refuse. We have implemented it, and this section reports
what the implementation established --- including three respects in which
building it corrected the specification, and one respect in which it
revealed a gap the specification did not know it had.

\begin{figure*}[!htbp]
\centering
\includegraphics[width=\textwidth]{figures/panel1-closure-dynamics.png}
\caption{\textbf{Closure decides the character of the dynamics, not merely
its magnitude.}
(\textbf{A})~State of a corticospinal circulation over $12$\,s, closed (left
axis) and after severance of the cortical output element (right axis). The
two are plotted on separate scales because they differ by four orders of
magnitude: the closed circulation executes bounded oscillation about a
non-zero centre, as \Cref{cor:perpetual} requires of an ungrounded graph
under sustained supply, while the severed one ramps monotonically until it
saturates the integrator. No value of a gain interpolates between these.
(\textbf{B})~Phase portrait of the closed circulation, state against its time
derivative over $9$\,s. The trajectory fills an annular region and never
approaches the marked centroid, which is the geometric content of
\Cref{thm:no-equilibrium}: the system has no fixed point to relax onto, so
what electrophysiology calls a resting potential is the centre of a cycle
rather than a state the circuit occupies.
(\textbf{C})~Excursion envelope on a logarithmic axis, closed and open, with
the median divergence time across all open arms in the corpus marked. The
closed envelope is stationary near $10^{-2}$; the open envelope crosses it
within the first second and saturates four decades higher.
(\textbf{D})~Envelope of the closed circulation as a function of loop gain
and elapsed time, with level contours projected onto the base plane. Decay
is governed by the product of gain and time, so a weakly conducting return
sustains a wider limit cycle without ever failing to close --- the surface
approaches zero envelope only asymptotically and never leaves the closed
regime.}
\label{fig:panel1}
\end{figure*}

\subsection{The reference implementation}

A complete implementation---lexer, LL(1) recursive-descent parser, type
checker, the four static analyses, the operational semantics of
\Cref{sec:semantics}, and a continuous backend discharging obligations
(B1)--(B4)---has been written and exercised against ten programs. It is a
\emph{reference} implementation in the strict sense: its purpose is to make
the specification's claims executable rather than merely stated, so that a
divergence between the two is a test failure rather than a matter of
opinion. Fifty tests assert specification properties directly. Among them:
\Cref{prop:attenuation} is checked at scaling factors down to $10^{-12}$;
\Cref{thm:lesion-monoid} is checked over all $3!$ orderings of three
distinct lesions; \Cref{thm:telescoping} is checked over $500$ randomly
generated cascades; and \Cref{prop:backend-independent} is checked by
confirming that the diagnostic stream is invariant under re-analysis.

The ten programs comprise $47$ experimental arms and yield $257$
observations, recorded with the provenance that (B2) and (B3) require: each
value carries the floor used to compute it, the frequency band it was
attributed to, the random seed, and the sample count.

\begin{figure*}[!htbp]
\centering
\includegraphics[width=\textwidth]{figures/panel2-attenuation-severance.png}
\caption{\textbf{Attenuation and severance are different operators, and the
difference is visible before any integration.}
(\textbf{A})~Closure index against the scaling factor $\lambda$ over twelve
decades. The index is invariant at \textsf{closed} throughout, confirming
\Cref{prop:attenuation} at factors down to $10^{-12}$; the marker at the
left edge is the value taken by the structurally distinct
\texttt{without element} operator. The discontinuity is not a limit of the
curve but a different point in the operator space.
(\textbf{B})~Tonic discharge rate over return gain and loop delay. Rate falls
smoothly with attenuation and rises as the reciprocal of delay, so a
parametric lesion moves the operating point continuously across this surface
while leaving the closure class fixed.
(\textbf{C})~Force output against outbound gain. The two curves --- one for
circulations that remain closed, one for those that have been opened ---
coincide exactly, because force is computed from the integrity and gain of
the outbound phase alone and closure is a property of both phases. The
markers are the measured nerve-block stages.
(\textbf{D})~The measured sequence, stages $1$--$6$ of the phased block, with
bar colour giving the closure index. Stage $3$ carries the full baseline
force of $420$\,N with the circulation open; stage $4$ carries none. Strength
and coordination therefore dissociate, and they dissociate discontinuously.}
\label{fig:panel2}
\end{figure*}

\subsection{Derivations that recover known clinical taxonomy}

Three programs reproduce classifications that clinical neurology arrived at
empirically. The interest is not in the conclusions, which are standard, but
in the derivation: none of these files encodes the conclusion anywhere. What
is declared is anatomy---compartments, delays, and two path declarations per
circulation---and the classification is the output of the linear-time
closure computation of \Cref{thm:closure-decidable}.

\paragraph{Upper versus lower motor neuron lesion.}
A program declares two circulations over shared compartments: a
corticospinal loop traversing cortex and brainstem, and a segmental stretch
reflex confined to the spinal level. Severing the cortical output element
opens the first; the second is not mentioned in that lesion, because the
lesion does not reach it. The observed result is

\begin{center}
\begin{tabular}{@{}lccl@{}}
\toprule
Lesion & Supraspinal loop & Segmental loop & Derived picture\\
\midrule
Cortical output & open & closed, $15.7$\,Hz & spasticity\\
Anterior horn   & open & open, no rhythm    & flaccidity\\
\bottomrule
\end{tabular}
\end{center}

The segmental loop survives a cortical lesion because it does not traverse
the severed element, and it continues to oscillate at the rate its own loop
delay implies---which is the circuit-level content of spasticity. The same
operator applied to the motor axon opens both, and the tonic rate becomes
undefined rather than small, because an open circulation sustains no rhythm
at all. This is the upper/lower motor neuron distinction obtained from
topology, and \Cref{cor:perpetual} is what makes the surviving oscillation a
prediction rather than an artefact.

\paragraph{Spinal cord injury.}
A second program declares circulations at three spinal levels, using the
template facility of \Cref{sec:extensions} to avoid restating the segmental
arc per muscle. A complete thoracic injury is expressed by severing
conduction through the thoracic level in both directions. Circulations
traversing that level open; the segmental arc below it remains closed at
$12.9$\,Hz; the upper-limb circulation, which traverses the cervical level
only, is untouched at $10.0$\,Hz. Paralysis below, preserved reflexes below,
and normal function above are three consequences of one membership test.

The program also expresses a motor-complete, sensory-incomplete injury by
applying \emph{different} operators to the descending and ascending
elements---severance to one, attenuation to the other. That this is
expressible at all is \Cref{prin:structural}; that the two receive different
closure indices is \Cref{prop:attenuation}.

\paragraph{Tremor.}
A third program applies four different lesion operators to one declared
postural circuit. All four leave the circulation closed, so the closure
index does not separate them and the spectral observables must. The measured
separation is

\begin{center}
\begin{tabular}{@{}lcccc@{}}
\toprule
Arm & reflex band & supraspinal band & coupling & $\eta$\\
\midrule
intact       & $0.241$ & $0.0123$ & $0.052$ & $0.899$\\
parkinsonian & $0.278$ & $0.0188$ & $0.048$ & $0.931$\\
essential    & $0.229$ & $0.0226$ & $0.063$ & $\mathbf{0.166}$\\
cerebellar   & $0.281$ & $0.0198$ & $0.047$ & $0.938$\\
\bottomrule
\end{tabular}
\end{center}

The final column is the reason \Cref{rem:report-eta} imposes a reporting
obligation, and it is worth dwelling on. The essential-tremor arm is the one
modelled by cross-stratum coupling, and it is the only arm whose type
separation collapses. This is not a defect of that arm's typing: it is
\Cref{prop:degeneracy} operating exactly as stated. Coupling across
non-adjacent strata disperses catalytic power \emph{within} each event type,
so the types stop separating, so any composed estimate from that arm is
uninterpretable. The language reports $\eta=0.166$ alongside $\kap=0.601$
rather than reporting the latter alone. A framework that returned only the
$\kap$ would have returned a confident number computed in precisely the
regime where confidence is unwarranted.

\begin{figure*}[!htbp]
\centering
\includegraphics[width=\textwidth]{figures/panel3-lesion-level.png}
\caption{\textbf{Which circulations survive a lesion is decided by which
traverse it.}
(\textbf{A})~Tonic rate of two circulations after severance of the cortical
output element. The supraspinal circulation, which traverses the severed
element, sustains no rhythm; the segmental circulation, which does not,
continues at $15.7$\,Hz. This pair is the circuit-level content of the
upper motor neuron picture, and neither value is declared anywhere in the
program.
(\textbf{B})~Fraction of circulations left open as a function of lesion level
for three anatomical classes: lower-limb supraspinal loops, upper-limb
supraspinal loops, and segmental arcs. The classes separate because a
membership test decides each: a loop opens exactly when the severed level
lies on one of its declared paths.
(\textbf{C})~Loop latency over segment index and declared per-hop delay. The
surface is planar because latency is a sum of declared delays, which is what
makes the reflex-latency budget of \Cref{def:strata} checkable against
measurement rather than fitted to it.
(\textbf{D})~Measured latencies of four circulations in the corpus, coloured
by closure index. The corticospinal loop is both the longest and the one the
cortical lesion opens.}
\label{fig:panel3}
\end{figure*}

\subsection{Findings not otherwise available}

Three programs produce statements that we have not found in the literature
and that follow from the closure formulation rather than from any parameter
choice.

\paragraph{Dissociation of strength from coordination under regional block.}
A phased program (\Cref{sec:extensions}) models the onset of a peripheral
nerve block in the order the fibre classes are recruited. The intermediate
stage---after the proprioceptive afferent has blocked but before the motor
efferent has---produces the following:

\begin{center}
\begin{tabular}{@{}lcc@{}}
\toprule
Stage & Closure index & Force output\\
\midrule
baseline            & closed        & $420.0$\,N\\
analgesia           & closed        & $420.0$\,N\\
proprioceptive loss & \textbf{open} & $\mathbf{420.0}$\,\textbf{N}\\
motor block         & open          & $0.0$\,N\\
motor recovery      & open          & $126.0$\,N\\
resolved            & closed        & $420.0$\,N\\
\bottomrule
\end{tabular}
\end{center}

Force is unchanged to full floating-point precision while the circulation is
open. The equality is not arranged: force is computed from the integrity and
gain of the outbound phase, closure from the topology of both phases, and
the two are independent functions of the declaration. Severing a return
element changes the second without touching the first.

The prediction is therefore that regional anaesthesia passes through a
window in which contractile capacity is fully preserved and coordinated
movement is nevertheless abolished---a dissociation of kind, not of degree.
This is the deafferentation phenomenology
\citep{cole1991,sainburg1995,ghez1995,lajoie1996} relocated into a
reversible, dose-controlled, routinely performed clinical procedure. An
optimal feedback control formulation, in which the afferent enters as a
multiplicative gain \citep{todorov2002}, predicts graded degradation across
this window instead. The two accounts disagree, the disagreement is sharp,
and the window is accessible.

\paragraph{Surgical repair as closure restoration.}
A program models trans-radial amputation, targeted muscle reinnervation, and
mirror therapy. Amputation is \texttt{without return}, which opens the
circulation. Reinnervation and mirror therapy are both \texttt{reroute},
which is the only operator in the language able to carry a circuit from
$\textsf{open}$ back to $\textsf{closed}$. The measured latencies are

\begin{center}
\begin{tabular}{@{}lcc@{}}
\toprule
Arm & Closure index & Loop latency\\
\midrule
intact    & closed & $49.8$\,ms\\
amputated & open   & ---\\
TMR       & closed & $65.8$\,ms\\
mirror    & closed & $35.8$\,ms\\
\bottomrule
\end{tabular}
\end{center}

Both interventions restore closure, and neither restores the original
circulation. The reinnervated loop is $16$\,ms slower than the limb it
replaces, because it traverses a different muscle with different conduction
distances; the mirror loop is faster but, as the checker notes at compile
time, traverses the supraspinal stratum, so it carries that stratum's
bandwidth and requires conscious mediation. The framework thus distinguishes
two repairs that both work, by the delay and the stratum of the substituted
path---a statement about prosthetic proprioception that the signal-and-gain
formulation cannot make, because in that formulation there is nothing to
restore, only a gain to raise.

\paragraph{Joint stiffness as an emergent quantity.}
A program declares two circulations sharing a spinal interneuron pool and a
joint, and measures their co-activation. Degrading reciprocal inhibition
raises the co-contraction ratio from $0.63$ to $0.87$ and the joint
stiffness from $139$ to $148$\,N/m. No stiffness is commanded anywhere in
the program; it is a property of two coupled limit cycles, and it moves
because the coupling changed. Under \Cref{cor:perpetual} this is what
stiffness must be: if the circuits admit no static equilibrium, a joint held
between them has no equilibrium stiffness either, only the effective
stiffness of two oscillations that overlap to a degree the coupling sets.

\begin{figure*}[!htbp]
\centering
\includegraphics[width=\textwidth]{figures/panel4-spectral-separation.png}
\caption{\textbf{When closure does not separate lesions, spectra must ---
and the separation statistic says whether they can.}
(\textbf{A})~Power spectral density of four tremor arms over $40$\,s. All
four circulations are closed, so the closure index is uninformative here and
the distinction is carried entirely by the distribution of power across the
stratum bands.
(\textbf{B})~Reflex-band and supraspinal-band power for the same four arms.
The parkinsonian and cerebellar arms elevate reflex-band power by
approximately one seventh; the essential arm does not, and instead raises
supraspinal-band power most.
(\textbf{C})~The diagnostic plane: type separation $\eta$ against coupling
index. Three arms sit above $0.89$; the essential arm, the only one modelled
by coupling across non-adjacent strata, falls to $0.166$, below the flagging
threshold. This is \Cref{prop:degeneracy} realised on real output --- the
cross-stratum edge disperses catalytic power within each event type, so the
typing stops separating and any composed estimate from that arm is
uninterpretable regardless of how confident it looks.
(\textbf{D})~Log spectral density over frequency and arm. The reflex-band
ridge is common to all four; what differs is its height relative to the
low-frequency shoulder, which is the quantity panel~\textbf{B} extracts.}
\label{fig:panel4}
\end{figure*}

\subsection{Prosthetic compliance as capacitance}

One result exercises \Cref{sec:compartments} rather than the closure
analysis. A program declares left and right hip circulations identical in
topology, delays, and gains, differing in exactly one declared quantity: the
terminal compartment's capacitance, $1.8\times10^{-4}$\,F for the biological
joint and $0.9\times10^{-4}$\,F for the prosthetic one. The measured forces
are $474.5$\,N and $335.6$\,N, a ratio of $0.7071$. This is
$\sqrt{0.9/1.8}$ to six figures, which is \eqref{eq:charge} evaluated at the
two capacitances.

The point is not that the implementation performs a square root correctly.
It is that a mechanical property of an implanted device---joint
compliance---enters the model as a compartment capacitance, propagates
through the charge relation without further modelling assumptions, and
emerges as a force asymmetry. Rule I is what makes this traceable: because
charge carries the compartment it was computed against, the two limbs'
charges are not interchangeable quantities and the asymmetry cannot be
silently averaged away.

\begin{figure*}[!htbp]
\centering
\includegraphics[width=\textwidth]{figures/panel5-repair-closure.png}
\caption{\textbf{Repair restores closure through a path that is not the
original.}
(\textbf{A})~Loop latency of the intact circulation, the amputated
circulation, and two repairs, with the intact value marked. Amputation
carries no latency because there is no loop to traverse.
(\textbf{B})~Latency offset of each repair from the original. Targeted muscle
reinnervation closes the circulation $16$\,ms slower, because the
substituted return traverses a different muscle at a different conduction
distance; mirror substitution closes it $14$\,ms faster but through the
supraspinal stratum, which the checker reports at compile time and which
carries that stratum's bandwidth.
(\textbf{C})~Restored latency over the number of hops in the substituted
return and the delay per hop, with level contours on the base plane. Any
point on this surface is a closed circulation; the surgical question is
which point, and the framework answers it before any integration.
(\textbf{D})~Closure index against loop latency for the four arms. Repair
moves a circulation vertically, from open to closed, and horizontally, to a
latency the original did not have. A signal-and-gain formulation has no
vertical axis here: there is nothing to restore, only a gain to raise.}
\label{fig:panel5}
\end{figure*}

\subsection{The degenerate estimator, exhibited}

\Cref{thm:telescoping} is a theorem and needs no empirical support. It is
nevertheless the claim in this paper that readers find least intuitive,
because the two quantities are computed by routines that look independent.
We therefore exhibit it. The implementation computes the instance-specific
prediction and the measured net power over $4000$ cascades from each of
three generating processes: one obeying the composition law, one constructed
to violate it, and one that is pure noise.

\begin{center}
\begin{tabular}{@{}lcc@{}}
\toprule
Generating process & $\max\abs{\hat\kap-\kap_{\mathrm{meas}}}$ & Pearson $r$\\
\midrule
multiplicative (obeys the law)  & $2.2\times10^{-16}$ & $1.000000$\\
adversarial (violates the law)  & $1.8\times10^{-15}$ & $1.000000$\\
random walk (no process at all) & $1.4\times10^{-15}$ & $1.000000$\\
\bottomrule
\end{tabular}
\end{center}

The agreement is at machine precision for data generated by a process
constructed to violate every substantive claim of the algebra, and for data
generated by no process whatsoever. A statistic that reports perfect
agreement on a random walk is not measuring the random walk.

Replacing the estimator with the typed one of \Cref{def:type-averaged} makes
the discrepancy nonzero, as \Cref{thm:nondegenerate} requires, and the
separation statistic behaves as \Cref{prop:degeneracy} predicts: a corpus
with well-separated types gives $\eta=0.934$, and one whose type means agree
to two figures gives $\eta=0.013$. On multiplicatively generated data with
mixed-type cascades the four laws of \Cref{prop:competitors} are cleanly
ordered, the multiplicative law attaining $\mathrm{RMSE}=0.033$ and
$r=0.967$ against $0.19$, $0.32$, and $0.44$ for its competitors---so the
corrected test selects, which is what a test with a null hypothesis is
supposed to do.

\begin{figure*}[!htbp]
\centering
\includegraphics[width=\textwidth]{figures/panel6-coupled-stiffness.png}
\caption{\textbf{Joint stiffness is emergent from two coupled circulations,
not commanded by either.}
(\textbf{A})~Activation of an agonist and an antagonist circulation sharing a
spinal interneuron pool and a joint, over $4$\,s with reciprocal inhibition
intact. The two alternate: reciprocal inhibition through the shared pool
pushes them out of phase, which is what smooth alternating movement is.
(\textbf{B})~The same pair after the inhibition is degraded by a declared
cross-stratum lesion. The traces now move together, because the mechanical
coupling through the shared joint is left unopposed.
(\textbf{C})~Measured co-contraction ratio and joint stiffness for the two
conditions. The ratio rises from $0.63$ to $0.87$ and the stiffness from
$139$ to $148$\,N/m. No stiffness is commanded anywhere in the program; it
moves because the coupling changed.
(\textbf{D})~Stiffness over reciprocal inhibition strength and the fraction
of each circulation's path that the shared junctions occupy. Stiffness falls
as inhibition rises and as sharing increases, and is defined only where the
two paths genuinely intersect --- a pair whose nominated junctions lie on
neither path is reported as uncoupled rather than silently coupled.}
\label{fig:panel6}
\end{figure*}

\subsection{What implementation changed in the specification}

Three corrections arose from building the system, and we record them because
each was a specification-level error that reading could not have caught.

\emph{An observable must name what computes it.} The specification treats
observables as opaque names discharged by a backend. Exercising the language
revealed that this admits names such as \texttt{conscious\_overhead}
alongside \texttt{loop\_latency}, with nothing to distinguish a quantity
that has a measurement procedure from one that does not. An observable
without a procedure is a promise, not a result. The implementation therefore
maintains a registry in which each admissible observable is declared with
its unit, its arity, and the procedure that computes it, and the checker
rejects any name absent from it. We regard this as a repair to
\Cref{sec:backends} rather than an implementation detail: it is the
difference between a framework that produces results and one that produces
vocabulary.

\emph{Force is a property of the outbound phase alone.} An early backend
computed force from the simulated state trace. This reported an open
circulation as producing more force than an intact one, because the open
circulation diverges and the divergence is large. The error inverted the
central claim of the nerve-block experiment. The correct treatment is that
the outbound phase carries the command to the muscle, so contractile
capacity depends on the integrity and gain of that phase and not on whether
the return phase closes---which is precisely why strength and coordination
dissociate. The specification did not say this because the specification did
not have to compute anything.

\emph{Loop gain is a property of the whole circulation.} A second backend
computed loop gain from the return path alone, which made lesions of
descending elements silent. Attenuating a descending pathway must change the
dynamics, since the restoring term traverses the entire loop. Both errors
share a shape: they arise from treating one phase of a circulation as
privileged, which is exactly the habit the language exists to break, and
they are the sort of error that survives careful reading and dies on contact
with a test.

\begin{figure*}[!htbp]
\centering
\includegraphics[width=\textwidth]{figures/panel7-degenerate-estimator.png}
\caption{\textbf{The instance-specific estimator agrees with its own
measurement on data generated by no process at all.}
(\textbf{A})~Composed prediction against directly measured net power for
cascades drawn from three generating processes: one obeying the composition
law, one constructed to violate it, and one a random walk. Every point of
every process lies on the identity line. If the two quantities were
independent, the adversarial and random-walk clouds would scatter.
(\textbf{B})~Maximum absolute discrepancy over $4000$ cascades per process,
on a logarithmic axis, with double-precision epsilon marked. All three sit
within one order of magnitude of it. A statistic that reports perfect
agreement on a random walk is not measuring the random walk
(\Cref{cor:no-null}).
(\textbf{C})~Residual over cascade length and process index. The surface is
flat in the process direction --- the agreement is insensitive to what
generated the data, which is the signature of an identity rather than a
confirmation --- and rises only with length, as accumulated floating-point
error.
(\textbf{D})~Root-mean-square error of four composition laws under the typed
estimator, on multiplicatively generated data with mixed-type cascades. The
multiplicative law attains $0.033$ against $0.19$, $0.32$, and $0.44$; the
corrected test therefore selects, which is what a test possessing a null
hypothesis is supposed to do.}
\label{fig:panel7}
\end{figure*}

\subsection{Extensions the experiments required}
\label{sec:extensions}

Exercising the language against clinically realistic programs revealed that
the core grammar of \Cref{sec:grammar} is expressive enough for ablation but
not for two other things experimenters need. We describe the extensions
and, equally, the ones we declined.

\emph{Templates.} A realistic model declares one segmental arc per muscle
group per level, differing only in compartment names and conduction delays.
A template facility, \texttt{circuit template} with instantiation, removes
the repetition. It is pure syntactic sugar: each instance expands to a
circuit declaration before typechecking, so every result of
\Cref{sec:typing,sec:semantics} applies unchanged to the expanded form.

\emph{Reroute.} \texttt{reroute return($v$) through $p$} replaces a
circulation's return phase with an explicit new path. It is the only
operator that can raise the closure index, and this asymmetry is the point:
without it the language models ablation but not repair, and reinnervation,
prosthetic feedback, and sensory substitution are all repairs. The aperture
analysis extends without modification, since the substituted path is checked
against the outbound phase exactly as the original was; a reroute that fails
to terminate at the outbound origin leaves the circuit open, and the
aperture report says why.

\emph{Phases.} The semantics of \Cref{sec:semantics} applies every lesion
before any observation, which makes any experiment whose circuit changes
over time inexpressible. A \texttt{phase} block, optionally naming a
predecessor, resolves this. Each phase is an independent experiment sharing
the intact circuit, so progress, preservation, and termination apply per
phase and the step bound becomes
$\sum_j(\abs{L_j}+\abs{O_j})+n_{\mathrm{ph}}$. Without phases the
nerve-block dissociation reported above cannot be posed at all, since its
content is entirely in the ordering.

\emph{Antagonist pairs.} Two circulations acting on one joint share
compartments. This is not composition of two models---which \Cref{sec:scope}
declines---but current conservation at a shared vertex within one medium,
and the distinction matters: the shared vertices must lie on both
circulations' declared paths, and the implementation reports a pair whose
nominated junctions do not as uncoupled, rather than silently coupling loops
that do not touch.

We declined two further extensions. \emph{Stochastic element failure} was
proposed to model fatigable weakness at the neuromuscular junction. It would
make the closure index a random variable, and hence a runtime rather than a
static property, forfeiting \Cref{prop:backend-independent}. It also buys
nothing: fatigable weakness is a trajectory through safety-factor loss,
which phases express while leaving closure statically decidable, and the
clinical distinction it was proposed to capture---myasthenic weakness
against denervation---is already the distinction between attenuation and
severance. Our implementation obtains it directly: the myasthenic arm is
closed with $63.0$\,N and an intact $12.9$\,Hz rhythm, the denervated arm
open with $0.0$\,N and no rhythm. \emph{Bilateral declaration} was proposed
for left--right comparison. It is metadata: every typing and reduction rule
applies to each circuit independently, so it would add a keyword without
adding semantics, and the comparison is properly performed over results.

\begin{figure*}[!htbp]
\centering
\includegraphics[width=\textwidth]{figures/panel8-capacitance-force.png}
\caption{\textbf{A mechanical property of an implant enters as a compartment
capacitance and emerges as a force asymmetry.}
(\textbf{A})~Force output against terminal compartment capacitance, with the
two declared values marked: $1.8\times10^{-4}$\,F for the biological hip and
$0.9\times10^{-4}$\,F for the prosthetic one. The curve is
\eqref{eq:charge}; the markers are measurements.
(\textbf{B})~Force over capacitance and outbound gain, with level contours on
the base plane. The two variables enter differently --- as a square root and
linearly --- so a compliance change and a weakness are not
interchangeable, which is what the compartment index of Rule~I preserves.
(\textbf{C})~Force ratio against capacitance ratio. The measured pair falls
on $\sqrt{0.5}=0.7071$ to six figures, which is \eqref{eq:charge} evaluated
at the two declared capacitances and not a fitted relation.
(\textbf{D})~Three observables of the prosthetic limb relative to the
biological one, with the $\sqrt{0.5}$ reference marked. Latency and
oscillation amplitude are unchanged, since topology, delays, and gains are
identical by construction; only force moves, and it moves to the reference
line. The asymmetry is attributable to one declared quantity.}
\label{fig:panel8}
\end{figure*}

\subsection{Threats to validity}
\label{sec:threats}

The backend used here is a linearised reference implementation. It
integrates a delayed multi-stratum loop and does not integrate cross-bridge
kinetics or musculotendon mechanics, so its absolute magnitudes---amplitudes
in arbitrary units, forces scaled from a nominal tetanic maximum---are not
physiological predictions and should not be read as such. What survives this
caveat is precisely what \Cref{prop:backend-independent} says survives: the
closure indices, the aperture reports, the compartment and stratum analyses,
and quantities derived from declared values by \eqref{eq:charge} are
backend-independent. The force ratio of $0.7071$ is exact; the $420$\,N it
scales is a convention.

Two further limits deserve statement. The findings above are consequences of
declarations, and a model may declare wrongly: the type system guarantees
that compartments, strata, and floors are used consistently, not that they
are true (\Cref{sec:scope}). And the clinical reproductions demonstrate that
the method recovers known taxonomy from anatomy, not that the anatomy was
chosen without knowledge of the taxonomy. We have tried to make that
inference checkable by keeping the declarations minimal and stating, for
each program, what is not encoded in it; a reader who suspects the
capacitances or delays were tuned to the answer can verify that no program
mentions spasticity, flaccidity, or any clinical category.

% =====================================================================
\section{Discussion}
\label{sec:discussion}
% =====================================================================

\subsection{Relation to motor control theories}

The closure account developed here is not a rival to existing motor-control
frameworks across their whole domain, but it differs from each on a specific
prediction, and the differences are what a Vitruvius experiment is designed to
expose.

The equilibrium-point hypothesis \citep{feldman1986} holds that a movement
command specifies a set point toward which the limb converges by the elastic
and reflexive properties of the musculotendon system. Within the present
framework the set point is recovered as the centre of the bounded oscillation
of a closed loop, and the width of that oscillation is the physiological
tremor observed in steady posture. The two accounts agree on convergence
without explicit feedback correction and differ on whether the target is a
point in configuration space or a cycle in state space; the second predicts
residual oscillation that the first must add separately.

Optimal feedback control \citep{todorov2002} treats the controller and the
plant as distinct systems joined by feedback whose gain enters
multiplicatively. It therefore predicts that removing feedback increases
variance while preserving the mean trajectory. The closure account predicts
instead a change of kind, because severance empties the constraint set of
\Cref{thm:min-dissipation} rather than perturbing its minimiser. The
deafferentation literature reports abolition rather than degradation of
coordinated movement \citep{cole1991,sainburg1995,ghez1995,lajoie1996}, which
is the pattern the closure account predicts and the gain formulation does not.

Internal-model theories \citep{kawato1999,wolpert2000} hold that the nervous
system maintains forward and inverse models updated by efference copy and
sensory return, and therefore predict that a deafferented subject can move
using internal prediction with errors accumulating as the model drifts. The
closure account requires no separate model: the loop's own return path
supplies the predictive dynamics, and learning is modification of circuit
conductance. The two differ on whether a deafferented subject should be able
to execute a well-practised movement open-loop for a short period, which is an
experimentally addressable question.

We state plainly that these are differences of prediction, not demonstrations.
Vitruvius is an instrument for posing them, and \Cref{lst:ablation} is the
shape such a posing takes.

\subsection{What the design buys}

The language makes four things checkable before any computation. A circulation
that has been opened is named, with the prediction that follows. A charge
computed at one capacitance cannot be silently combined with one computed at
another. An influence between non-adjacent time scales must pass through the
intervening scale and incur its delay, unless it is declared as a lesion. And
a floor of zero, which would make the fractional observables undefined, cannot
be declared. Because none of these consults a backend
(\Cref{prop:backend-independent}), they are available at the cost of parsing.

It also removes a specific way of being wrong. The instance-specific estimator
of \Cref{def:instance} is the natural thing to write, is implemented by
routines that look independent, and agrees with its own prediction to machine
precision on every data set. \Cref{thm:telescoping} shows the agreement is an
identity; rule IV makes the identity unwritable.

\subsection{What it does not buy}

The language does not make a model correct. Compartment indices, strata, and
floors are declarations, and a model may declare them wrongly; what the type
system guarantees is that the declarations are used consistently, not that
they are true. The stratum discipline enforces adjacency but cannot verify
that a given element belongs to the stratum it claims. The typed estimator has
a genuine null hypothesis but no power when type separation is low
(\Cref{prop:degeneracy}), and the language can only report $\eta$, not repair
a typing that fails to carve the system.

Nor does the closure analysis certify physiological realism. A closed circuit
is one whose declared paths match; whether those paths correspond to
anatomical pathways is outside what the compiler can see. \Cref{lst:reflex}
would typecheck with the delays permuted.

Finally, the design is deliberately narrow. It has no representation for
prescribed kinematics, externally specified controllers, or open-loop
actuation of an intact system, all of which are legitimate objects of study in
general-purpose platforms \citep{delp2007,seth2018,todorov2012}. A user who
needs them needs a different tool, or needs to treat the Vitruvius model as
one component in a larger workflow.

% =====================================================================
\section{Scope and limitations}
\label{sec:scope}
% =====================================================================

We record the limitations that a reader should weigh before adopting the
language.

\emph{The lumped representation.} \Cref{sec:foundations} treats compartments
as points at which potential is uniform. Spatially extended phenomena---cable
propagation along an axon, diffusion in the synaptic cleft, continuum
mechanics of muscle---are represented by lumped delays and conductances. This
is adequate for establishing closure and its consequences and inadequate for
detailed waveform prediction, which requires the partial differential
formulations that a continuous backend supplies internally.

\emph{Strata are a coarse partition.} \Cref{def:strata} recognises three
levels. Slower processes exist, including drift over minutes, and faster ones,
including physiological tremor above $5$\,Hz. A model requiring these must
either extend the stratum lattice---which the adjacency rule of
\Cref{def:stratum-rule} accommodates without modification, being defined over
an arbitrary linear order---or accept their absence.

\emph{Closure is topological, not functional.} The analysis of
\Cref{sec:aperture} decides whether declared paths match. It does not decide
whether a return path carries enough information to be useful, and it will
report a circuit closed when the return path is intact but its conductance has
been scaled to a value at which the loop is functionally silent. The
distinction between an attenuated loop and a severed one is exactly what
\Cref{prop:attenuation} is about, and users should note that the language
takes the topological side of it by design.

\emph{The floor is only as good as its derivation.} \Cref{prop:floor}
guarantees positivity for a connected closed circuit, and
\Cref{prop:floor-sensitivity} quantifies the cost of misestimation, but
neither validates a particular \texttt{derived} expression. A model that
discharges the floor by a sample minimum satisfies the type system while
carrying no evidential weight (\Cref{rem:floor-estimators}).

\emph{Backends are trusted.} Obligations (B1)--(B4) are requirements on a
backend, not properties the language verifies. A backend that violates (B1) by
raising on divergence converts the central result of an ablation experiment
into a crash, and nothing in the type system prevents it.

\emph{No cross-model composition.} A program describes one system. Composing a
musculoskeletal model with models of other physiology is outside the language,
by design: what would be required is an account of how the compartments of two
models share a medium, and we do not supply one here.

% =====================================================================
\section{Conclusion}
\label{sec:conclusion}
% =====================================================================

Vitruvius has one primitive, four non-standard typing rules, and three
reduction rules. The primitive is the circulation, and it is chosen because
the intact neuromuscular graph has no external ground
(\Cref{prop:no-ground}), therefore no static equilibrium
(\Cref{thm:no-equilibrium}), and therefore no way for an efferent projection
to terminate except by closing (\Cref{cor:closure-obligatory}). From that
premise the language's distinguishing capability follows: severing a return
pathway and attenuating it are different operators with different closure
indices, distinguishable statically (\Cref{prop:attenuation}), where a
signal-and-gain formulation makes them the same limit.

Three further analyses fall out of the same structure. Charge carries the
compartment against which it was computed, so mixing capacitances is a type
error rather than a silent numerical mistake. Influence between time scales
must traverse the intervening scale, so a cross-stratum shortcut is either a
declared lesion or a rejected program. And a fractional observable requires a
positive floor, whose derivation the source must state, because an
upward-biased floor inflates exactly the quantities describing events nearest
the irreducible limit.

The metatheory is standard---progress, preservation, order-independent lesion
composition, termination---and the backend interface is narrow enough that the
same program runs unchanged over integrators of quite different character,
with the static diagnostics unaffected by the choice
(\Cref{prop:backend-independent}).

The result we regard as most useful is negative. The natural estimator of an
event's contribution agrees with its own prediction on every data set that can
exist (\Cref{thm:telescoping}), so a perfect correlation between a composed
prediction and its measurement is not a confirmation of anything. The language
makes that estimator unwritable and requires the typed alternative, which can
fail, together with the separation statistic without which its success cannot
be interpreted. A language for experiments should make the vacuous experiment
hard to express, and this is the one place where we have been able to do so
mechanically rather than by advice.

\bibliographystyle{plainnat}
\bibliography{references}

\end{document}
